% AUTOGENERATED DRAFT from manuscript source: manuscript/chapters/04-particle-pushers.md
% Converter: scripts/convert_markdown_chapter_to_tex.py
% Review gate: docs/latex-migration-plan.md (Phase 2 per-chapter gate)

\chapter{粒子推进器：从 Lorentz 方程到 \texttt{PushPX()}}
\label{chap:04-particle-pushers}

粒子推进器负责把单个宏粒子从一个时间层推进到下一个时间层。它表面上是局部单粒子算法，实际上依赖上一章建立的主循环条件：场必须已经填好 gather guard cells，粒子位置和动量必须处在正确 leapfrog 时间层，外场、电离电荷态、辐射反作用和 QED 选项也必须在进入 pusher 前准备好。

本章按一条读者可追踪的因果链展开：先用 Lorentz 方程固定时间层和动量记号；再比较 Boris、Vay 与 Higuera--Cary 如何更新动量；最后从 \cpath{WarpX::PushParticlesandDeposit()} 经 \cpath{MultiParticleContainer::Evolve()} 进入 \cpath{PhysicalParticleContainer::Evolve()} 与 \cpath{PushPX()}，检查 gather、外场、push、位置更新和沉积如何接在同一 tile loop 中。需要核对实现时，优先搜索 \cpath{PushSelector.H} 的 \cpath{doParticleMomentumPush()}、三个 \cpath{UpdateMomentum*.H} 函数，以及 \cpath{PhysicalParticleContainer::PushPX()}；不要把行号或一次源码快照当作算法语义。

阅读 Boris 推进时要特别区分半步磁旋转的 Birdsall--Langdon 半角关系，不能把旋转系数机械地除以二；\cpath{Examples/Tests/particle\_pusher} 提供 Higuera--Cary force-free 路径的直接验证入口。

\subsection{阅读路线：先定位一条带电粒子的时间链}

第一次阅读不必逐个记住后面的多物理分支。先用以下四步建立一条能够排错的主线：

\begin{enumerate}
\item \textbf{先读 4.1--4.4。} 固定 $\mathbf{x}^n$、$\mathbf{u}^{n-1/2}$、$\mathbf{u}^{n+1/2}$ 的时间层，理解 Boris、Vay 与 Higuera--Cary 各自要解决的离散更新问题；这一步回答“更新的对象是什么”。
\item \textbf{再读 4.5--4.10。} 沿 \cpath{PushParticlesandDeposit()} 到 \cpath{PushPX()} 追踪一次带质量粒子的 tile loop：gather、外场叠加、动量更新、位置更新和沉积的顺序不能互换；这一步回答“这些公式何时真正发生”。
\item \textbf{按需要进入 4.11--4.14。} 辐射反作用、隐式推进、ionization、collisions、边界粒子与 QED 都会改变粒子状态或创建/移除粒子。选择其中一节时，先写清 source、product、时间位置和要比较的 observable，避免把独立的物理模型误当作 Boris 的一个参数。
\item \textbf{最后用 4.15--4.16 收束。} 把异常轨道、粒子数、场残差或守恒量映射回时间层、tile 主链和相应 analysis；练习要求的结论必须注明一个不能由该案例外推的范围。
\end{enumerate}

这条路线把本章的细节压缩成一个判断顺序：先问粒子状态在哪个时间层，再问哪条更新链消费它，最后才问某个扩展模型或回归能证明什么。

\section{连续 Lorentz 方程}

WarpX 对带质量粒子推进的基本方程是相对论 Lorentz 方程。令

\[
\mathbf{u}=\gamma\mathbf{v},
\qquad
\gamma=\sqrt{1+\frac{|\mathbf{u}|^2}{c^2}},
\qquad
\mathbf{v}=\frac{\mathbf{u}}{\gamma}.
\]

这里 WarpX 源码中的 \code{ux, uy, uz} 是 $\mathbf{u}$ 的三个分量，而不是普通速度 $\mathbf{v}$。对质量 $m$、电荷 $q$ 的粒子，

\[
\frac{d\mathbf{u}}{dt}
=
\frac{q}{m}\left(\mathbf{E}+\mathbf{v}\times\mathbf{B}\right),
\]

\[
\frac{d\mathbf{x}}{dt}
=
\mathbf{v}
=
\frac{\mathbf{u}}{\gamma}.
\]

显式 leapfrog PIC 中，常用时间层是

\[
\mathbf{x}^{n}\rightarrow \mathbf{x}^{n+1},
\qquad
\mathbf{u}^{n-1/2}\rightarrow \mathbf{u}^{n+1/2}.
\]

因此位置推进使用更新后的半步动量：

\[
\mathbf{x}^{n+1}
=
\mathbf{x}^{n}
+\frac{\mathbf{u}^{n+1/2}}{\gamma^{n+1/2}}\Delta t.
\]

WarpX 的显式位置更新位于 \cpath{Source/Particles/Pusher/UpdatePosition.H} 的 \cpath{UpdatePosition()}；它先通过 \cpath{GetExplicitPusherDisplacement()} 用完整三速度分量构造位移，再按编译几何写回实际存储的坐标分量。

这里可以补一个 Dawson 1983 的历史边界。那篇综述在 \code{electromagnetic particle models and fractional dimensional models} 一节明确指出：只要问题涉及自洽磁场或 electromagnetic radiation，particle model 就不能再被理解成“electrostatic 外面多加几个场分量”，而必须回到 full Maxwell equations 与 self-consistent current/charge representation。与此同时，低空间维度并不意味着低速度维度；为了在一维或二维空间里承载 electromagnetic wave，仍必须保留 transverse \cpath{E/B/j} 与 transverse velocity，于是才会自然出现 \code{1 1/2-D}、\code{1 2/2-D} 和 \code{2 1/2-D} 这些 fully electromagnetic reduced-dimension models。这个边界正好解释了为什么本章后面很多看似“低维”的 laser、beam 和 FEL benchmark，仍然必须认真讨论 magnetic rotation、relativistic momentum 与 transverse field coupling，而不能把它们误降成纯 electrostatic pusher。

\section{Boris 推进的结构}

Boris pusher 把动量更新拆成三个操作：

\begin{enumerate}
\item 电场半步；
\item 磁场旋转；
\item 电场半步。
\end{enumerate}

设旧动量为 $\mathbf{u}^{n-1/2}$，先做电半步：

\[
\mathbf{u}^{-}
=
\mathbf{u}^{n-1/2}
+\frac{q\Delta t}{2m}\mathbf{E}^{n}.
\]

再用 $\mathbf{u}^{-}$ 计算

\[
\gamma^-=\sqrt{1+\frac{|\mathbf{u}^-|^2}{c^2}},
\qquad
\mathbf{t}
=
\frac{q\Delta t}{2m\gamma^-}\mathbf{B}^{n},
\qquad
\mathbf{s}
=
\frac{2\mathbf{t}}{1+|\mathbf{t}|^2}.
\]

磁场旋转写成

\[
\mathbf{u}'=\mathbf{u}^-+\mathbf{u}^-\times\mathbf{t},
\]

\[
\mathbf{u}^{+}=\mathbf{u}^-+\mathbf{u}'\times\mathbf{s}.
\]

最后做第二个电半步：

\[
\mathbf{u}^{n+1/2}
=
\mathbf{u}^{+}
+\frac{q\Delta t}{2m}\mathbf{E}^{n}.
\]

\sourceline{源码文件}{\cpath{Source/Particles/Pusher/UpdateMomentumBoris.H}}
\sourceline{函数}{\cpath{UpdateMomentumBoris()}}

下面是为阅读重排的核心节选；变量名和控制流与源码一致，完整函数仍以该文件和函数名为准：

% <-- fenced code (cpp)
\begin{codeblock}
const amrex::ParticleReal econst = 0.5_prt*q*dt/m;
const bool first_or_full =
    momentum_push_type == MomentumPushType::FirstHalf ||
    momentum_push_type == MomentumPushType::Full;
const bool second_or_full =
    momentum_push_type == MomentumPushType::SecondHalf ||
    momentum_push_type == MomentumPushType::Full;
const bool split_half =
    momentum_push_type == MomentumPushType::FirstHalf ||
    momentum_push_type == MomentumPushType::SecondHalf;

if (first_or_full) {
    // First half-push for E
    ux += econst*Ex;
    uy += econst*Ey;
    uz += econst*Ez;
}
// Compute temporary gamma factor
constexpr auto inv_c2 = PhysConst::inv_c2_v< amrex::ParticleReal>;
const amrex::ParticleReal inv_gamma = 1._prt/std::sqrt(1._prt + (ux*ux + uy*uy + uz*uz)*inv_c2);
// Magnetic rotation
amrex::ParticleReal tx = econst*inv_gamma*Bx;
amrex::ParticleReal ty = econst*inv_gamma*By;
amrex::ParticleReal tz = econst*inv_gamma*Bz;
if (split_half) {
    const amrex::ParticleReal tsq = tx*tx + ty*ty + tz*tz;
    const amrex::ParticleReal factor = (tsq > 0._prt)
        ? (std::sqrt(1._prt + tsq) - 1._prt) / tsq : 0.5_prt;
    tx *= factor;
    ty *= factor;
    tz *= factor;
}
const amrex::ParticleReal tsqi = 2._prt/(1._prt + tx*tx + ty*ty + tz*tz);
const amrex::ParticleReal sx = tx*tsqi;
const amrex::ParticleReal sy = ty*tsqi;
const amrex::ParticleReal sz = tz*tsqi;
const amrex::ParticleReal ux_p = ux + uy*tz - uz*ty;
const amrex::ParticleReal uy_p = uy + uz*tx - ux*tz;
const amrex::ParticleReal uz_p = uz + ux*ty - uy*tx;
// - Update momentum
ux += uy_p*sz - uz_p*sy;
uy += uz_p*sx - ux_p*sz;
uz += ux_p*sy - uy_p*sx;
if (second_or_full) {
    // Second half-push for E
    ux += econst*Ex;
    uy += econst*Ey;
    uz += econst*Ez;
}
\end{codeblock}

% <-- table block (source: 04-particle-pushers L176)
\begin{tabularx}{\textwidth}{LLL}
\toprule
更新阶段 & 代码动作 & 公式对应 \\
\midrule
\endfirsthead
\toprule
更新阶段 & 代码动作 & 公式对应 \\
\midrule
\endhead
电场系数 & \code{econst = 0.5*q*dt/m} & $\frac{q\Delta t}{2m}$ \\
第一半步 & FirstHalf 或 Full 时做电半步 & $\mathbf{u}^{n-1/2}\to\mathbf{u}^{-}$ \\
旋转参数 & 计算 \cpath{inv\_gamma} 和 full-step \cpath{t} & $\gamma^{-},\mathbf{t}$ \\
分裂修正 & FirstHalf/SecondHalf 时按半角关系重标定 \cpath{t} & 使两次 half push 合成一次 Full 的磁旋转 \\
磁旋转 & 交叉乘更新动量 & $\mathbf{u}^{-}\to\mathbf{u}^{+}$ \\
第二半步 & SecondHalf 或 Full 时做第二个电半步 & $\mathbf{u}^{+}\to\mathbf{u}^{n+1/2}$ \\
\bottomrule
\end{tabularx}

函数注释说明 \cpath{FirstHalf} 和 \cpath{SecondHalf} 可以分裂执行，并且连续执行应与一次 \cpath{Full} 更新数学等价。这正好服务于 \cpath{WarpXEvolve.cpp} 中把碰撞放在 momentum push 中间的路径。

这里有一个必须区分的实现细节：WarpX 的 half momentum push 不是简单把 $q\,\Delta t\,B/(2m\gamma)$ 再乘 $1/2$。\cpath{UpdateMomentumBoris()} 的 half-push 分支使用

\[
\frac{|t_{\mathrm{half}}|}{|t_{\mathrm{full}}|}
=
\frac{\sqrt{1+|t_{\mathrm{full}}|^2}-1}{|t_{\mathrm{full}}|^2}
\]

重标定 \code{tx,ty,tz}。这来自 \cpath{tan(alpha/2)} 与 \cpath{tan(alpha/4)} 的半角关系，目的是让 \cpath{FirstHalf} 后接 \cpath{SecondHalf} 的组合仍对应 full Boris rotation，而不是两个朴素半磁旋转的近似拼接。

\section{WarpX 如何选择 pusher}

\sourceline{源码文件}{\cpath{Source/Particles/Pusher/PushSelector.H}}
\sourceline{函数}{\cpath{doParticleMomentumPush()}}

% <-- table block (source: 04-particle-pushers L202)
\begin{tabularx}{\textwidth}{LL}
\toprule
条件 & 选择 \\
\midrule
\endfirsthead
\toprule
条件 & 选择 \\
\midrule
\endhead
进入函数后 & species 电荷乘 \cpath{ion\_lev}，得到当前粒子的有效电荷。 \\
启用 classical radiation reaction & 进入 Boris 家族分支：通常调用 \cpath{UpdateMomentumBorisWithRadiationReaction()}；若编译 QED 且同步分支中的 $\chi$ 不低于门限，则调用普通 \cpath{UpdateMomentumBoris()}。 \\
\cpath{ParticlePusherAlgo::Boris} & 进入 \cpath{UpdateMomentumBoris()}。 \\
\cpath{ParticlePusherAlgo::Vay} & 进入 \cpath{UpdateMomentumVay()}。 \\
\cpath{ParticlePusherAlgo::HigueraCary} & 进入 \cpath{UpdateMomentumHigueraCary()}。 \\
\bottomrule
\end{tabularx}

这说明输入参数选择的不是一个高层“粒子模块”，而是每个粒子在 \cpath{PushPX()} device loop 内调用的单粒子 momentum update。下面继续展开 Vay 与 Higuera-Cary 的源码。

这一节背后的经典来源可以直接回到 Birdsall-Langdon 1985 第一分卷 \cpath{4-3} 到 \cpath{4-5}。那里把磁推进的核心先写成几何分裂：电场部分是半步 impulse，磁场部分是速度空间旋转；随后再把旋转压成 $t=\tan(\theta/2)$、$s=2t/(1+t^2)$、$c=(1-t^2)/(1+t^2)$ 这组半角变量，并进一步给出向量 Boris 形式。对本章来说，这个来源有两个价值。第一，WarpX 的 Boris 更新不是孤立经验公式，而是这条 “half-accel + rotation + half-accel” 离散合同的现代实现。第二，Birdsall 在 \cpath{4-5} 里明确区分了 \cpath{1d2v/1d3v} 和真正的一维动力学，这正好解释了为什么即使空间维数较低，本章后面讨论的 mover 仍必须保留多速度分量与磁旋转结构。

Boris 1970 的原始历史位置需要单独标注边界：本书给出 J. P. Boris 的会议论文书目和 DTIC \cpath{ADA023511} 入口，但目前没有可逐页核对的会议论文全文。因此，本章的算法推导采用 Birdsall--Langdon 1985 的完整讲解，WarpX 的实现说明则回到 \cpath{Source/Particles/Pusher/UpdateMomentumBoris.H}；Boris 1970 不能作为已经逐页核查的直接公式依据。

物理上可以先记住：

\begin{itemize}
\item Boris：经典、鲁棒、磁场部分近似旋转，长期性质好。
\item Vay：针对相对论漂移和 Lorentz 变换一致性问题设计，在 boosted frame 场景常用。
\item Higuera-Cary：相对论粒子推进的结构保持改进，常用于减少高相对论问题中的系统误差。
\end{itemize}

正式误差比较必须回到对应论文和 benchmark；本章先把 WarpX 的实际实现讲清楚。

\section{Vay pusher：相对论速度变换一致性的更新}

\sourceline{源码文件}{\cpath{Source/Particles/Pusher/UpdateMomentumVay.H}}
\sourceline{函数}{\cpath{UpdateMomentumVay()}}

源码注释明确引用 Vay 2008 的公式 (9)-(13)，并说明 \cpath{FirstHalf} 与 \cpath{SecondHalf} 连续执行应等价于 \cpath{Full}。下面是为阅读压缩的核心节选；它省略函数签名和重复分量，但保留时间层分支、$\gamma$ 根和旋转结构：

% <-- fenced code (cpp)
\begin{codeblock}
const amrex::ParticleReal econst = q*dt/m *
    ((momentum_push_type == MomentumPushType::Full) ? 1.0_prt : 0.5_prt);
const amrex::ParticleReal bconst = 0.5_prt*q*dt/m;
const amrex::ParticleReal inv_gamma = 1._prt/std::sqrt(
    1._prt + (ux*ux + uy*uy + uz*uz)*PhysConst::inv_c2_v<amrex::ParticleReal>);
    // Get tau
const amrex::ParticleReal taux = bconst*Bx;
const amrex::ParticleReal tauy = bconst*By;
const amrex::ParticleReal tauz = bconst*Bz;
const amrex::ParticleReal tausq = taux*taux + tauy*tauy + tauz*tauz;
const amrex::ParticleReal uxpr = ux + econst*Ex +
    ((momentum_push_type == MomentumPushType::SecondHalf)
        ? 0.0_prt : (uy*tauz - uz*tauy)*inv_gamma);
const amrex::ParticleReal uypr = uy + econst*Ey +
    ((momentum_push_type == MomentumPushType::SecondHalf)
        ? 0.0_prt : (uz*taux - ux*tauz)*inv_gamma);
const amrex::ParticleReal uzpr = uz + econst*Ez +
    ((momentum_push_type == MomentumPushType::SecondHalf)
        ? 0.0_prt : (ux*tauy - uy*taux)*inv_gamma);

if (momentum_push_type != MomentumPushType::FirstHalf) {
        // Get gamma'^2
    const amrex::ParticleReal gprsq = 1._prt +
        (uxpr*uxpr + uypr*uypr + uzpr*uzpr)*PhysConst::inv_c2_v<amrex::ParticleReal>;
        // Get u*
    const amrex::ParticleReal ust =
        (uxpr*taux + uypr*tauy + uzpr*tauz)*PhysConst::inv_c_v<amrex::ParticleReal>;
        // Get new gamma
    const amrex::ParticleReal sigma = gprsq - tausq;
    const amrex::ParticleReal gisq = 2._prt/(sigma + std::sqrt(
        sigma*sigma + 4._prt*(tausq + ust*ust)));
        // Get t, s
    const amrex::ParticleReal bg = bconst*std::sqrt(gisq);
    const amrex::ParticleReal tx = bg*Bx;
    const amrex::ParticleReal ty = bg*By;
    const amrex::ParticleReal tz = bg*Bz;
    const amrex::ParticleReal s = 1._prt/(1._prt + tausq*gisq);
        // Get t.u'
    const amrex::ParticleReal tu = tx*uxpr + ty*uypr + tz*uzpr;
        // Get new U
    ux = s*(uxpr + tx*tu + uypr*tz - uzpr*ty);
    uy = s*(uypr + ty*tu + uzpr*tx - uxpr*tz);
    uz = s*(uzpr + tz*tu + uxpr*ty - uypr*tx);
} else {
    ux = uxpr;
    uy = uypr;
    uz = uzpr;
}
\end{codeblock}

Vay pusher 和 Boris 的主要差别在求解更新后的相对论因子这一步。Boris 使用旧的 $\gamma^-$ 构造磁旋转；Vay 先构造 $\mathbf{u}'$，再通过

\[
\sigma=\gamma'^2-\tau^2,
\qquad
\gamma_\mathrm{new}^{-2}
=
\frac{2}{\sigma+\sqrt{\sigma^2+4(\tau^2+u_*^2)}}
\]

解析得到新的 $\gamma$ 相关量。源码中的 \cpath{gisq} 就是这里的 $\gamma_\mathrm{new}^{-2}$，\cpath{bg=bconst*sqrt(gisq)} 对应 $(q\Delta t/2m)/\gamma_\mathrm{new}$。这使磁旋转使用与更新后状态一致的相对论因子，适合 boosted-frame 或高速束流穿越问题。

\cpath{FirstHalf} 路径只返回 \code{u'}，\cpath{SecondHalf} 路径不会再加入旧速度磁项。这和 Boris 文件中的分裂设计一致，服务于碰撞、外力或隐式/显式混合路径。

把这段源码重新放回 Vay 2008 原文，边界会更清楚。Vay 论文的第一性问题不是“再发明一个 relativistic mover”，而是对 relativistic crossing / boosted-frame 场景，离散 Lorentz force 必须尽量保住 electric 和 magnetic contributions 的 cancellation property。作者先用

\[
\mathbf E+\mathbf v\times\mathbf B=0
\]

这个试金石说明 Boris 在一般非零 \cpath{E}、\cpath{B} 情况下会出现 spurious force，然后才提出新的 leapfrog velocity average。于是 WarpX 的 \cpath{UpdateMomentumVay.H} 最值得保留的历史定位不是“Boris 的另一个变体”，而是：它实现的是一条以 frame-change consistency 为目标的专门修正路线。

Vay 2008 的 \cpath{II.C} 单粒子测试不是普通轨道展示，而是同一物理系统在 laboratory frame 和 moving frame 下都要和解析解一致的 frame-consistency test。常量 $B_z$ 例子以 $v_x=10^{-2}c$ 起步，并取 $\Delta t=10^{-2}\times 2\pi/\omega_c$：新 pusher 在实验室系和沿 $\hat y$ 方向、$\gamma_f=2$ 的 moving frame 中都贴住解析轨道；Boris 即使加上 $\tan(\omega_c\Delta t)/(\omega_c\Delta t)$ 修正，也会在 moving frame 中偏离，且误差在 $\gamma_f=3$ 后迅速放大。

常量 $E_x=1\,\mathrm{kV/m}$ 的测试给出相同结论：电子在实验室系初始静止，经历 \cpath{100} 步、每步 $1\,\mathrm{ns}$ 的更新时，三种 mover 在实验室系里都正确；只有新 pusher 在 $\gamma_f=100$ 的 moving frame 中仍保持解析一致。因此这篇文献的硬证据不是“Vay 在某些 case 里更稳”，而是 Boris 的误差会在 frame change 后由次要项变成主导项。

这篇论文也没有顺手给出一个通用 Maxwell solver。它在 \cpath{III} 节明确把场求解边界限定在 waves 和 retardation 可忽略、并且对每个 species 可以在共动系里近似取

\[
v_z \gg v_x,v_y,
\qquad
\frac{\partial}{\partial t}\approx v_z\frac{\partial}{\partial z}
\]

的场景。于是 field side 被压成带 $\gamma z$ 拉伸的 Poisson 型求解，并近似保留 electrostatic、magnetostatic 以及沿主流向的 inductive effect。对 \cpath{N} 个 species，代价就是 \cpath{N} 次这类 Poisson solve。这条有边界的 Darwin-lite explicit approximation 解释了为什么 \cpath{IV} 节的 LHC-like ultrarelativistic beam / electron-cloud 应用会特意选在 $\gamma\approx16.5$ 的 moving frame 中做 first-principles PIC：在那里 beam 与 electron cloud 的 self-electric / self-magnetic cancellation 最强，最能放大 mover 的 frame-consistency 缺陷。

文中报告 Boris 无论是否带 $\tan$ 修正，都会让 beam 和 electron 宏粒子以非物理速度丢失；只有新 pusher 才能恢复预期的 hose-like instability，并给出与实验室系 quasistatic WARP calculation 一致的 vertical emittance growth rate 和 saturation level。因此，对 WarpX 而言，\cpath{UpdateMomentumVay.H} 的历史角色应理解成 relativistic beam-crossing / boosted-frame consistency repair，而不是一个与一般场求解器或一般 relativistic mover 等价并列的“备选算法”。

\subsection{Vay Appendix A/B：显式 \texorpdfstring{$\gamma$}{γ} 根与回旋半径边界}

Vay 2008 的 Appendix A 给出了源码中 \cpath{gisq} / \cpath{gamma\_new} 公式为什么可以显式计算。磁旋转中先定义

\[
\mathbf u^{i+1}=s\left[\mathbf u' +(\mathbf u'\cdot\mathbf t)\mathbf t+\mathbf u'\times\mathbf t\right],
\qquad s=\frac{1}{1+t^2},
\qquad \mathbf t=\frac{\boldsymbol\tau}{\gamma^{i+1}},
\]

其中 $\boldsymbol\tau=(q\Delta t/2m)\mathbf B$。对上式与 $\mathbf u$ 做点积，利用 $\gamma^2=1+u^2/c^2$，并令

\[
\gamma'=\sqrt{1+u'^2/c^2},
\qquad u^*=\frac{\mathbf u'\cdot\boldsymbol\tau}{c},
\qquad \sigma=\gamma'^2-\tau^2,
\]

可把隐式的 relativistic factor 压成一个关于 $\gamma^2$ 的二次方程：

\[
\gamma^4+(\tau^2-\gamma'^2)\gamma^2-\tau^2-u^{*2}=0.
\]

只保留正的实根，得到

\[
\gamma^{i+1}=\sqrt{\frac{\sigma+\sqrt{\sigma^2+4(\tau^2+u^{*2})}}{2}}.
\]

这解释了 \cpath{UpdateMomentumVay.H} 的实现顺序：先用 \code{u'} 和 \cpath{tau} 构造标量不变量，再取正根，最后由 $t=\tau/\gamma$ 和 $s=1/(1+t^2)$ 完成旋转。\cpath{gisq} 存的是 $\gamma^{-2}$，因此 device loop 不需要迭代求解 $\gamma$。这是 Appendix A 与当前 kernel 的直接公式桥接，不是普通 Boris 旋转中的经验系数。

Appendix B 则给出常磁场、$\mathbf E=0$ 时的 gyroradius 边界：

\[
\Delta\theta=2\arctan\left(\frac{\omega_c\Delta t}{2}\right),
\qquad
\mathbf v^{i+1/2}=\left[1+\left(\frac{\omega_c\Delta t}{2}\right)^2\right]\frac{\mathbf v^i+\mathbf v^{i+1}}{2},
\]

从圆周几何得到

\[
R=\frac{\|\mathbf v^{i+1/2}\|}{\omega_c}
=\left[1+\left(\frac{\omega_c\Delta t}{2}\right)^2\right]^{1/2}\frac{\|\mathbf v^i\|}{\omega_c}.
\]

所以“Vay 在任意时间步都给出正确 gyroradius”必须加限定：若用于位置推进的半步速度满足 $\|\mathbf v^{i+1/2}\|=v_0$，则 $R=v_0/\omega_c$；若把整数时刻速度直接当作 $v_0$，仍会出现与 Boris 类似的放大因子。pusher 的动量更新、半步速度定义和位置更新必须一起检查。

一个有用的最小验证是去掉自洽场、AMR 和 PML，只保留均匀磁场，然后同时比较三类量：离散相位、由相邻位置差构造的速度 proxy、以及由该 proxy 得到的 gyroradius。这样的比较能够检查“公式中的半步速度是否与实际位置更新相容”，也能把 Higuera--Cary 的相位行为与 Boris/Vay 分开观察；它不能证明输出文件中存在可直接读取的 half-step 速度，更不能代替论文图形的逐点复现。读者应把 Appendix B 当成一个关于时间层的判别题，而不是把任何圆形轨道都视作该附录结论的证明。

Vay 2008 与 Higuera--Cary 2017 在本章承担不同作用：前者说明为何需要对相对论速度变换作修正，后者说明如何在 Boris-like 结构中构造相对论不变量。两篇论文都应回到原文公式和本章给出的 kernel 对照阅读；只有在输入、观察量和容差明确相同的情况下，才可以讨论某个 WarpX 例子是否复现了其中的专门结论。

\subsection{用推进器谱系约束 Vay 的结论范围}

Vay--Godfrey 2014 review 的读者价值，不是替 WarpX 的 \cpath{UpdateMomentumVay.H} 背书，而是把 Boris、Lorentz-invariant pusher、场更新、current deposition、field gather、filtering 与数值稳定性放在同一条 PIC 离散链上。它提醒读者：推进行为不能只凭一条单粒子轨迹判断，场与源项怎样被离散、怎样被 gather，同样会决定相对论计算的误差结构。

因此第 4 章应按四层证据阅读：Vay 2008 解释 frame-consistency 这一原始算法目标；该综述给出推进器在完整 PIC 方法谱系中的位置；WarpX 源码说明 kernel 实际消费的变量和时间层；明确输入和观察量的算例才说明某个条件下实际测到了什么。任何一层都不能替代另一层，尤其不能把综述中的历史算法图或其他 PIC 程序的结果写成 WarpX 的验证结论。

本节的核心判断不依赖于资料整理方式：Vay 是为特定的相对论 frame-consistency 问题设计的推进器。选择它之前，仍要同时检查粒子推进、场更新、沉积和诊断路径，而不能只依据一个 mover 名称或一条单粒子轨道。

\section{Higuera-Cary pusher：Boris-like 结构的相对论修正}

\sourceline{源码文件}{\cpath{Source/Particles/Pusher/UpdateMomentumHigueraCary.H}}
\sourceline{函数}{\cpath{UpdateMomentumHigueraCary()}}

% <-- fenced code (cpp)
\begin{codeblock}
template <typename T>
AMREX_GPU_HOST_DEVICE AMREX_INLINE
void UpdateMomentumHigueraCary(
    T& ux, T& uy, T& uz,
    const T Ex, const T Ey, const T Ez,
    const T Bx, const T By, const T Bz,
    const T q, const T m, const amrex::Real dt )
{
    using namespace amrex::literals;

    // Constants
    const T qmt = 0.5_prt*q*dt/m;
    // Compute u_minus
    const T umx = ux + qmt*Ex;
    const T umy = uy + qmt*Ey;
    const T umz = uz + qmt*Ez;
    // Compute gamma squared of u_minus
    T gamma = 1._prt + (umx*umx + umy*umy + umz*umz)*PhysConst::inv_c2_v<T>;
    // Compute beta and betam squared
    const T betax = qmt*Bx;
    const T betay = qmt*By;
    const T betaz = qmt*Bz;
    const T betam = betax*betax + betay*betay + betaz*betaz;
    // Compute sigma
    const T sigma = gamma - betam;
    // Get u*
    const T ust = (umx*betax + umy*betay + umz*betaz)*PhysConst::inv_c_v<T>;
    // Get new gamma inverse
    gamma = 1._prt/std::sqrt(0.5_prt*(sigma + std::sqrt(sigma*sigma + 4._prt*(betam + ust*ust)) ));
    // Compute t
    const T tx = gamma*betax;
    const T ty = gamma*betay;
    const T tz = gamma*betaz;
    // Compute s
    const T s = 1._prt/(1._prt+(tx*tx + ty*ty + tz*tz));
    // Compute um dot t
    const T umt = umx*tx + umy*ty + umz*tz;
    // Compute u_plus
    const T upx = s*( umx + umt*tx + umy*tz - umz*ty );
    const T upy = s*( umy + umt*ty + umz*tx - umx*tz );
    const T upz = s*( umz + umt*tz + umx*ty - umy*tx );
    // Get new u
    ux = upx + qmt*Ex + upy*tz - upz*ty;
    uy = upy + qmt*Ey + upz*tx - upx*tz;
    uz = upz + qmt*Ez + upx*ty - upy*tx;
}
\end{codeblock}

前半段仍是电场半步：

\[
\mathbf{u}^-=\mathbf{u}^{n-1/2}+\frac{q\Delta t}{2m}\mathbf{E}.
\]

随后源码用 \cpath{beta=qmt*B} 和 \cpath{u\_minus} 计算

\[
\sigma = \gamma_-^2-\beta^2,
\qquad
u_* = \frac{\mathbf{u}^-\cdot\boldsymbol{\beta}}{c},
\]

并通过平方根表达式得到新的 \cpath{gamma}，这里变量名 \cpath{gamma} 在这一行之后实际保存的是 $\gamma^{-1}$。\code{tx,ty,tz} 是 $\boldsymbol{\beta}/\gamma$，$s=1/(1+t^2)$。\cpath{u\_plus} 的形式像 Boris 的磁旋转，但最后不是单纯加第二个电半步，而是

\[
\mathbf{u}^{n+1/2}
=
\mathbf{u}^+
+\frac{q\Delta t}{2m}\mathbf{E}
+\mathbf{u}^+\times\mathbf{t}.
\]

这个额外叉乘项是 Higuera-Cary 结构和 Boris 结构最容易混淆的地方。源码中 \cpath{upy*tz-upz*ty}、\cpath{upz*tx-upx*tz}、\cpath{upx*ty-upy*tx} 正是 $\mathbf{u}^+\times\mathbf{t}$ 的三个分量。

把这段 kernel 放回 Higuera-Cary 2017 原文，WarpX 里这条算法线的真实边界会更清楚。那篇论文并不是在 \code{Vay 2008} 的 boosted-frame cancellation 问题上继续竞争，而是把 Boris、Vay 和新方法放到三个并列判据下比较：\cpath{E=0} 时的能量守恒、crossed \cpath{E/B} 场下正确的 $\mathbf E\times\mathbf B$ drift、以及 phase-space volume preservation。作者的核心判断是：Boris 保住 volume 但 drift 不对；Vay 保住 drift 但不 volume-preserving；Higuera-Cary 则是三者中唯一同时保住 volume 与 $\mathbf E\times\mathbf B$ drift 的二阶 relativistic momentum integrator。因此，\cpath{UpdateMomentumHigueraCary.H} 最准确的历史定位不是“又一个 Boris-like 变体”，而是：在 Boris 的 rotation skeleton 上，把 centered average 和 $\gamma$ 的 prescription 改写成一条双结构保持路线。

Higuera-Cary 2017 的 \cpath{III-VI} 节还把这个判断压成了 WarpX 读者真正需要的两层证据。第一层是实现证据：新方法表面上仍是 implicit centered scheme，但最终可以像 Boris/Vay 一样显式实现，真正的实现分叉点几乎全部浓缩在 $\gamma_{new}$ 的求法上。这正好解释了为什么 WarpX 源码里 \cpath{UpdateMomentumHigueraCary.H} 的外形与 Boris 非常接近，却在 \cpath{sigma}、\cpath{ust} 和新的 relativistic factor 路径上分叉。第二层是数值/几何证据：作者用 Poincare surface of section 而不是普通能量曲线来比较 practical timestep 下的轨道拓扑。小时间步时三种方法都能给出嵌套曲线；但在更接近实际模拟的 $\Delta t=1/10$ 下，Vay 会出现 resonance island 和不同轨道 section 交叉，而 Boris 与 Higuera-Cary 仍保持正确的 phase-space topology。对本章来说，这意味着 Higuera-Cary 的价值判断标准不是 \code{Vay 2008} 那种 frame-change consistency，而是 geometric/topological preservation at practical timestep。

如果进一步把论文公式和 WarpX kernel 逐式对位，这条关系会更直观。Higuera-Cary 论文先定义

\[
\vec{\epsilon}=\frac{q\Delta t}{2m}\vec E,
\qquad
\vec{\beta}=\frac{q\Delta t}{2m}\vec B,
\qquad
\vec u_-=\vec u_i+\vec\epsilon,
\]

而源码里的 \cpath{qmt}、\cpath{umx/umy/umz}、\cpath{betax/betay/betaz} 正是这些量。随后论文显式求解

\[
\gamma_{new}^2
=
\frac{1}{2}
\left(
\gamma_-^2-\beta^2
+\sqrt{(\gamma_-^2-\beta^2)^2+4(\beta^2+|\vec\beta\cdot\vec u_-|^2)}
\right),
\]

WarpX 则不先存 $\gamma_{new}^2$，而是直接用 \code{sigma = gamma - betam}、$u_*=(\mathbf u_-\cdot\boldsymbol\beta)/c$ 和下一行的平方根公式，把变量 \cpath{gamma} 改写成 $\gamma_{new}^{-1}$。这一点很关键：代码里的 \cpath{gamma} 在函数中段被重载了，前半段是 $\gamma_-^2$，后半段则变成 rotation 要消费的 inverse relativistic factor。之后 \cpath{tx/ty/tz}、\cpath{s}、\cpath{umt}、\cpath{upx/upy/upz} 这条链，就是论文 Boris-like rotation equation

\[
\vec u_+ - \vec u_- = (\vec u_+ + \vec u_-) \times \frac{\vec\beta}{\gamma_{new}}
\]

的直接实现。因此，\cpath{UpdateMomentumHigueraCary.H} 最本质的实现差异可以压成一句话：它在 Boris 的 rotation skeleton 上，仅通过改写 $\gamma$ prescription，就把 \cpath{volume-preserving} 与 $\mathbf E\times\mathbf B$ drift preservation 这两条性质同时保住了。

论文第 V 节的 Jacobian 证明也让这一点不再停留在口头层。作者证明新 integrator 的前后半步 Jacobian determinant 互为倒数，所以一步更新的总体 Jacobian 恰好等于 \cpath{1}；而 Vay 的 Jacobian 一般写成 \code{J(x_i,u_i)/J(x_i,u_f)} 这样的比值，通常不会化成 \cpath{1}。这正是后面数值例子里 Vay 在 practical timestep 下出现 resonance island 和轨道交叉，而 Higuera-Cary 仍保持 phase-space topology 的几何根源。

如果把这段证明再往下压一层，最关键的中间对象其实是 $I-\Omega$。Higuera-Cary 论文先把后半步对 $\bar u_{new}$ 的 Jacobian 写成 “Boris-like rotation 主干 $I-\Omega$ + 一个 rank-one correction”，其中 $\Omega\cdot V=(\beta\times V)/\gamma_{new}$。这一步意味着作者不是直接对整个复杂映射硬算 determinant，而是先把旋转骨架和 relativistic correction 拆开。随后利用 determinant lemma，把后半步体积变化写成

\[
J_{f,new}=\det(I-\Omega)\times(\text{scalar correction}),
\]

再经过代数整理压成

\[
J _ { f , new } = 1 + \frac { \beta ^ { 2 } + ( \vec { \beta } \cdot \bar { u } _ { new } ) ^ { 2 } } { \gamma _ { new } ^ { 4 } } .
\]

前半步可得同一形式的 determinant，因此真正的 volume-preserving 不是“每个子步都单独等于 1”，而是前后半步 Jacobian determinant 互为逆，整步更新的 Jacobian 才严格等于 \cpath{1}。这一点很重要，因为它把 \cpath{UpdateMomentumHigueraCary.H} 的稳定性来源从经验判断提升成了明确的 Jacobian 结构断言。

这里还要额外提醒一个记号陷阱：论文里 $J_{f,new}$ 和 $J_{i,new}$ 最后都被压成相同的显式标量函数，但这不表示它们是“同一个 Jacobian”。它们对应的是后半步和前半步那两条相反方向映射上的 determinant，因此恰恰是因为它们处在 reciprocal 位置，整步 Jacobian 才会严格回到 \cpath{1}。同样，论文对 Vay 的结论也不是“任何情况下都会立刻出现 attractor/repeller”。作者保留了一个例外边界：若磁场在时空上恒定，\code{J(x_i,u_i)/J(x_i,u_f)} 这串比值会 telescoping，再结合 \code{J(x,u)} 在有界区域里的有界性，不能直接推出灾难性体积失真。真正的问题是它缺少一般性的 volume-preservation，因此在 practical timestep 和更复杂轨道拓扑下更容易暴露出 resonance-island 与 trajectory-crossing 这类非物理后果。

这组 regression 和文献的配对仍需保守表述。\cpath{Examples/Tests/particle\_pusher} 提供 force-free Higuera-Cary 强断言；Poincare 合同则验证 \code{x=0,p_x>0} 截面、\cpath{I\_y} 顺序和解析 quartic reference。topology classifier 同时保留时间顺序和相空间中心角顺序；在 32³、2201-frame 长轨道上，后三种 pusher 的角排序候选均无自交或轨道间交叉，说明原先时间折线的交叉计数是连接顺序伪影，而不是物理 resonance-island 证据。14-species dense family 与 64³ \cpath{p\_y=1.6/1.8} control 进一步显示 Vay 窗口漂移约 \cpath{6.5e-2}，控制组约 \cpath{1e-3}，但该 resonance-sensitive screen 仍不是 two-fold island 或 trajectory-crossing topology proof。

这些短轨道、长轨道、密集 \cpath{p\_y} family、resonance screen 和 resolution screen 给出的证据等级应当分开理解：短轨道采样不足；长轨道的 invariant/reference 与 angular-order candidate 可以通过，但 topology 仍需要人工复核；密集族的 resonance-sensitive screen 可以通过，但解析 reference curve 和跨 pusher 的特征尚未形成完整拓扑证明。因此，最强可写结论是“invariant 与局部 resonance-sensitive screen 已建立，论文等价的 topology gate 尚未启用”。

\section{从 \texttt{MultiParticleContainer} 到 \texttt{PhysicalParticleContainer}}

\sourceline{源码文件}{\cpath{Source/Evolve/WarpXEvolve.cpp}}
\sourceline{函数}{\cpath{WarpX::PushParticlesandDeposit()}}

它选择 current 字段名后调用 \cpath{mypc->Evolve(...)}。

\cpath{mypc} 是 \cpath{MultiParticleContainer}。

\begin{itemize}
\item 源码文件：\cpath{Source/Particles/MultiParticleContainer.cpp}
\item 函数：\cpath{MultiParticleContainer::Evolve()}
\end{itemize}

% <-- table block (source: 04-particle-pushers L529)
\begin{tabularx}{\textwidth}{LL}
\toprule
调度阶段 & 操作 \\
\midrule
\endfirsthead
\toprule
调度阶段 & 操作 \\
\midrule
\endhead
开始 & 不跳过沉积时清零 \cpath{current\_fp/current\_buf/rho\_fp/rho\_buf}。 \\
implicit 分支 & 处理隐式 solver 相关源项和 mass matrix 的额外清零逻辑。 \\
容器遍历 & 调用每个 \cpath{pc->Evolve(...)}。 \\
\bottomrule
\end{tabularx}

这层只负责多物种调度。真正的单 species 粒子推进在 \cpath{PhysicalParticleContainer::Evolve()}。

但在继续进入 \cpath{PushPX()} 之前，还要先看清容器层次和属性系统，否则后面很多变量名会被误读。

\cpath{Source/Particles/WarpXParticleContainer.H} 中的 \cpath{PIdx} 定义编译期属性表。它规定每个粒子天生就有：

\begin{itemize}
\item 位置分量 \cpath{x/y/z} 或非笛卡尔等价量；
\item 权重 \cpath{w}；
\item proper velocity \cpath{ux/uy/uz}；
\item 在 RZ/球坐标几何下额外的 \cpath{theta/phi}。
\end{itemize}

而 \cpath{IntIdx::nattribs} 默认是 0，这意味着 WarpX 的整数粒子属性默认都不是编译期内建，而是后续按需动态添加。

顶层类层次由 \cpath{Source/Particles/MultiParticleContainer.cpp} 与各容器头文件共同决定：

% <-- fenced code (text)
\begin{codeblock}
MultiParticleContainer
  -> WarpXParticleContainer
       -> PhysicalParticleContainer
            -> PhotonParticleContainer
            -> RigidInjectedParticleContainer
       -> LaserParticleContainer
\end{codeblock}

这里几类容器的物理职责不同：

\begin{itemize}
\item \cpath{WarpXParticleContainer} 是统一基类，提供粒子 SoA 骨架、gather/deposition/push 的共用接口；
\item \cpath{PhysicalParticleContainer} 承担普通带质量 species 的主要物理路径；
\item \cpath{PhotonParticleContainer} 继承 \cpath{PhysicalParticleContainer}，但其 \cpath{DepositCharge()} 和 \cpath{DepositCurrent()} 直接空实现，因此它保留很多粒子基础设施，却不承担带电沉积；定义见 \cpath{Source/Particles/PhotonParticleContainer.H}；
\item \cpath{LaserParticleContainer} 则直接从 \cpath{WarpXParticleContainer} 继承，因为激光天线粒子只需要 prescribed motion 和 current deposition，不需要普通 \cpath{FieldGather}；定义见 \cpath{Source/Particles/LaserParticleContainer.H}。
\end{itemize}

这层类分工直接决定了粒子属性的来源。\cpath{Source/Particles/PhysicalParticleContainer.cpp} 中的 \cpath{PhysicalParticleContainer} 构造函数会按模块开关注册第一批 runtime attributes：

\begin{itemize}
\item QED quantum synchrotron 时加 \cpath{opticalDepthQSR}；
\item Breit-Wheeler 时加 \cpath{opticalDepthBW}；
\item \cpath{addRealAttributes} / \cpath{addIntegerAttributes} 时加入用户 parser 驱动属性；
\item \cpath{save\_previous\_position} 时加 \cpath{prev\_x/prev\_y/prev\_z}。
\end{itemize}

电离模块的 \cpath{InitIonizationModule()} 稍后又会动态补上 \cpath{ionizationLevel}。因此 WarpX 的粒子属性系统不是一个静态结构体，而是：

\begin{enumerate}
\item 编译期 builtin real：\cpath{x/y/z/w/ux/uy/uz/...}；
\item 构造期或模块初始化期动态加入的持久物理状态：如 \cpath{opticalDepthQSR}、\cpath{opticalDepthBW}、\cpath{prev\_*}、\cpath{ionizationLevel}；
\item 更晚才按算法路径加入的临时缓存属性。
\end{enumerate}

最典型的临时属性来自 implicit solver：

\begin{itemize}
\item 函数：\cpath{ImplicitSolver::CreateParticleAttributes()}
\item 源码：\cpath{Source/FieldSolver/ImplicitSolvers/ImplicitSolver.cpp}
\end{itemize}

它会统一给粒子容器补加：

\begin{itemize}
\item \cpath{x\_n/y\_n/z\_n}
\item \cpath{ux\_n/uy\_n/uz\_n}
\item 如果启用 particle suborbits，再加 \cpath{nsuborbits}
\end{itemize}

而且都用 \code{comm = 0} 注册，所以这些量既不参与通信，也不写入 checkpoint。step-start 初始化的来源为：

\begin{itemize}
\item 源码：\cpath{Source/FieldSolver/ImplicitSolvers/WarpXImplicitOps.cpp}
\item 动作：把当前位置和动量快照写入 \cpath{x\_n} 和 \cpath{ux\_n}，并把 \cpath{nsuborbits} 先置成 1。
\end{itemize}

它们的角色不是长期物理属性，而是 implicit 时间推进器本步用的局部状态。

\cpath{WarpXParticleContainer::AddNParticles()} 进一步说明了这套系统怎样落地。

\begin{itemize}
\item 源码：\cpath{Source/Particles/WarpXParticleContainer.cpp}
\item 顺序：先写 builtin \cpath{x/y/z/w/ux/uy/uz}，再写调用者显式提供的 runtime real/int，最后调用 \cpath{DefaultInitializeRuntimeAttributes()} 给剩余 runtime attrs 自动补默认值。
\end{itemize}

于是：

\begin{itemize}
\item \cpath{opticalDepthQSR/BW} 可以由 QED engine 随机初始化；
\item \cpath{ionizationLevel} 可以统一设成 \cpath{ionization\_initial\_level}；
\item 用户 parser 属性可以按 \code{attribute.<name>(x,y,z,ux,uy,uz,t)} 自动求值。
\end{itemize}

也就是说，进入 \cpath{PushPX()} 之前，WarpX 已经把“粒子是什么物种、当前带哪些附加物理状态、哪些只是临时算法缓存”这三件事分清了。后面看到 \cpath{ux}、\cpath{ux\_n}、\cpath{ionizationLevel}、\cpath{opticalDepthQSR} 这些名字时，必须先回到这里判断它们属于哪一层状态。

\section{\texttt{PhysicalParticleContainer::Evolve()} 的 tile loop}

本章最重要的入口是 \cpath{Source/Particles/PhysicalParticleContainer.cpp} 中的 \cpath{PhysicalParticleContainer::Evolve()}。它把一个 species 的粒子按 tile 遍历，并把沉积、gather、push、buffer、隐式路径和 load-balance cost 放在一个局部循环里。

核心顺序是：

% <-- table block (source: 04-particle-pushers L616)
\begin{tabularx}{\textwidth}{LLL}
\toprule
tile-loop 阶段 & 操作 & 含义 \\
\midrule
\endfirsthead
\toprule
tile-loop 阶段 & 操作 & 含义 \\
\midrule
\endhead
准备 & 取得 \cpath{Efield\_aux} 和 \cpath{Bfield\_aux} & 粒子 gather 使用 auxiliary fields。 \\
条件判断 & 判断是否沉积 charge/current、是否 split particles & \cpath{skip\_deposition} 和 \cpath{do\_not\_deposit} 会关掉沉积。 \\
AMR 分区 & 遍历 tile，必要时按 AMR buffer 分区粒子 & fine/coarse gather 和 deposit 的粒子集合可能不同。 \\
push 前 & 沉积 \cpath{rho} component 0 & 旧时间层电荷，通常对应 $\rho^n$。 \\
fine patch & 粒子调用 \cpath{PushPX()} & gather fine fields 并推进粒子。 \\
buffer/coarse & 粒子调用 \cpath{PushPX()} & AMR 边界附近可从 coarse auxiliary fields gather。 \\
push 后 & 沉积 current & 显式路径 \cpath{relative\_time=-0.5*dt}，对应 $\mathbf{J}^{n+1/2}$。 \\
新时间层 & 沉积 \cpath{rho} component 1 & 新时间层电荷，通常对应 $\rho^{n+1}$。 \\
可选后处理 & particle splitting & subcycling 时避免 coarse level 重复沉积。 \\
\bottomrule
\end{tabularx}

这段源码说明，真实粒子推进不是“先推所有粒子，再单独沉积”。WarpX 为了 AMR、缓存局部性、GPU/CPU 并行和时间层一致性，在 tile 内完成 gather/push/deposit 的组合。

\section{\texttt{PushPX()}：gather 和 push 的融合 kernel}

\sourceline{源码文件}{\cpath{Source/Particles/PhysicalParticleContainer.cpp}}
\sourceline{函数}{\cpath{PhysicalParticleContainer::PushPX()}}

它是真正进入单粒子并行循环的地方。

核心源码原文如下，省略了 QED 宏分支和部分属性准备：

% <-- fenced code (cpp)
\begin{codeblock}
amrex::ParallelFor(
    TypeList<CompileTimeOptions<no_exteb,has_exteb>, CompileTimeOptions<no_qed  ,has_qed>>{},
    {exteb_runtime_flag, qed_runtime_flag},
    np_to_push,
    [=] AMREX_GPU_DEVICE (long ip, auto exteb_control, auto qed_control)
{
    amrex::ParticleReal xp, yp, zp;
    getPosition(ip, xp, yp, zp);

    amrex::ParticleReal Exp = Ex_external_particle;
    amrex::ParticleReal Eyp = Ey_external_particle;
    amrex::ParticleReal Ezp = Ez_external_particle;
    amrex::ParticleReal Bxp = Bx_external_particle;
    amrex::ParticleReal Byp = By_external_particle;
    amrex::ParticleReal Bzp = Bz_external_particle;

    if (gather_fields) {
        // first gather E and B to the particle positions
        doGatherShapeN(xp, yp, zp, Exp, Eyp, Ezp, Bxp, Byp, Bzp,
                       ex_arr, ey_arr, ez_arr, bx_arr, by_arr, bz_arr,
                       ex_type, ey_type, ez_type, bx_type, by_type, bz_type,
                       dinv, xyzmin, lo, n_rz_azimuthal_modes,
                       nox, galerkin_interpolation);
    }

    if constexpr (exteb_control == has_exteb) {
        getExternalEB(ip, Exp, Eyp, Ezp, Bxp, Byp, Bzp);
    }

    scaleFields(xp, yp, zp, Exp, Eyp, Ezp, Bxp, Byp, Bzp);

    doParticleMomentumPush<0>(ux[ip], uy[ip], uz[ip],
                                  Exp, Eyp, Ezp, Bxp, Byp, Bzp,
                                  ion_lev ? ion_lev[ip] : 1,
                                  mass, q, pusher_algo, do_crr,
                                  dt, momentum_push_type);

    if (position_push_type == PositionPushType::Full) {
        UpdatePosition(xp, yp, zp, ux[ip], uy[ip], uz[ip], dt, mass);
        setPosition(ip, xp, yp, zp);
    }
});
\end{codeblock}

关键步骤：

% <-- table block (source: 04-particle-pushers L686)
\begin{tabularx}{\textwidth}{LL}
\toprule
kernel 阶段 & 操作 \\
\midrule
\endfirsthead
\toprule
kernel 阶段 & 操作 \\
\midrule
\endhead
入口准备 & 检查 gather level、构造 gather box，并按 \cpath{ngEB} 扩展 guard cells。 \\
属性准备 & 准备粒子位置访问器、外场、动量数组、ionization level、旧位置缓存，以及 pusher/RR/QED 选项。 \\
并行分派 & 启动带 compile-time option 的 \cpath{amrex::ParallelFor}。 \\
场构造 & 读取粒子位置，调用 \cpath{doGatherShapeN()}；随后叠加每粒子外场 callback，再应用 \cpath{scaleFields()}。 \\
动量更新 & 调用 \cpath{doParticleMomentumPush()} 更新动量。 \\
位置更新 & 若 \cpath{PositionPushType::Full}，调用 \cpath{UpdatePosition()} 更新位置。 \\
\bottomrule
\end{tabularx}

这给出了 WarpX 显式粒子推进的真实顺序：

% <-- fenced code (text)
\begin{codeblock}
for particle in tile:
    read x^n and u^(n-1/2)
    gather E/B at x^n
    add external particle fields and apply field scaling
    push momentum to u^(n+1/2)
    push position to x^(n+1)
\end{codeblock}

随后 \cpath{PhysicalParticleContainer::Evolve()} 沉积半步电流与新时间层电荷。

\textbf{读者的单粒子状态检查卡。} 追踪一颗宏粒子时，要把“决定轨道的量”和“决定它对网格 source 的量”分开记账：

\begin{enumerate}
\item 进入 kernel 的动力学状态是 $\mathbf{x}^n$、$\mathbf{u}^{n-1/2}$、物种质量 $m$ 与有效电荷 $q_{\mathrm{eff}}=q\,\texttt{ionizationLevel}$。\cpath{doParticleMomentumPush()} 接收的是这个 $q_{\mathrm{eff}}$ 和 $m$，并不接收宏粒子权重 $w$。
\item pusher 实际消费的场不是某个原始 \cpath{MultiFab} 的单个值。常量 particle external field 先写入局部 \cpath{Exp...Bzp}；\cpath{doGatherShapeN()} 将 \cpath{Efield\_aux/Bfield\_aux} 的插值\textbf{累加}到这些局部量；每粒子 external-field functor 再继续叠加，最后 \cpath{scaleFields()} 才给出 $\mathbf{E}_{\mathrm{push}},\mathbf{B}_{\mathrm{push}}$。因此轨道诊断必须说明它比较的是网格场、外场，还是二者合成后的 pusher 场。
\item 只有 \cpath{PositionPushType::Full} 才会把更新后的 $\mathbf{u}^{n+1/2}$ 交给 \cpath{UpdatePosition()} 并写回 $\mathbf{x}^{n+1}$。一次 \cpath{PushPX()} 调用、一次 momentum half push 与一次完整物理位移不是同义词。
\item 随后的 charge/current deposition 才构造 $wq=q\,w\,\texttt{ionizationLevel}$。因此两个宏粒子即使具有相同的 $\mathbf{x},\mathbf{u},q/m$ 而权重不同，也会有相同的单粒子轨道、不同的 $\rho/\mathbf J$ 贡献。
\end{enumerate}

这张卡直接限定验证结论：单粒子轨道、force-free pusher 或 Larmor 类案例能够检查 gather/push/position 的局部合同；它们不能单独证明宏粒子 source、连续性或自洽场演化正确。后者必须连同权重、沉积算法和 \cpath{SyncCurrentAndRho()} 进入第 5 章的验证链。

\section{\texttt{doGatherShapeN()}：从网格场到粒子场}

\cpath{PushPX()} 在调用 pusher 前先调用 \cpath{doGatherShapeN()}。这一步的物理含义是把网格上的 Yee/PSATD 场插值到粒子位置：

\[
\mathbf{E}_p
=
\sum_{\mathbf{i}} \mathbf{E}_{\mathbf{i}} S_{\mathbf{i}}(\mathbf{x}_p),
\qquad
\mathbf{B}_p
=
\sum_{\mathbf{i}} \mathbf{B}_{\mathbf{i}} S_{\mathbf{i}}(\mathbf{x}_p).
\]

但 WarpX 不能只用一个标量形函数，因为 $E_x,E_y,E_z,B_x,B_y,B_z$ 在交错网格上的中心位置不同；同时 Galerkin 插值会让某些分量使用低一阶形函数。运行时入口在 \cpath{Source/Particles/Gather/FieldGather.H}。其完整函数签名包含粒子坐标、六个输出分量、六个场数组、六个 \cpath{IndexType} 和几何元数据；对理解分派而言，下面是等价的阅读伪代码：

% <-- fenced code (cpp)
\begin{codeblock}
// args 表示完整签名中的坐标、场数组、centering 与几何元数据。
if (galerkin_interpolation) {
    if (nox == 1) doGatherShapeN<1,1>(args);
    if (nox == 2) doGatherShapeN<2,1>(args);
    if (nox == 3) doGatherShapeN<3,1>(args);
    if (nox == 4) doGatherShapeN<4,1>(args);
} else {
    if (nox == 1) doGatherShapeN<1,0>(args);
    if (nox == 2) doGatherShapeN<2,0>(args);
    if (nox == 3) doGatherShapeN<3,0>(args);
    if (nox == 4) doGatherShapeN<4,0>(args);
}
\end{codeblock}

这里的 \cpath{nox} 不是运行时循环里的 shape 阶数变量，而是被转成模板参数 \cpath{depos\_order}。这样 GPU kernel 内部可以用 \code{if constexpr} 展开阶数，避免每个粒子再做阶数分支。\cpath{galerkin\_interpolation} 同理变成第二个模板参数，后面直接影响数组长度 \code{depos_order + 1 - galerkin_interpolation}。

模板主体开头也在 \cpath{Source/Particles/Gather/FieldGather.H}。下面是 x 方向有效计算的阅读伪代码；完整模板签名以及与 x 方向同构的 y/z 参数从略。\cpath{NODE}/\cpath{CELL} 的判断仍是源码中的判断：

% <-- fenced code (cpp)
\begin{codeblock}
constexpr int NODE = amrex::IndexType::NODE;
constexpr int CELL = amrex::IndexType::CELL;
Compute_shape_factor<depos_order> const shape;
Compute_shape_factor<depos_order - galerkin_interpolation> const gshape;
const amrex::Real x = (xp - xyzmin.x)*dinv.x;

if (component_needs_node) j_node = shape(sx_node, x);
if (component_needs_cell) j_cell = shape(sx_cell, x - 0.5_rt);
if (component_needs_node_galerkin) j_node_v = gshape(sx_node_galerkin, x);
if (component_needs_cell_galerkin) j_cell_v =
    gshape(sx_cell_galerkin, x - 0.5_rt);

// ex/by/bz 选择 Galerkin 权重；其余分量选择常规权重。
const auto& sx_ex = (ex_type[0] == NODE) ? sx_node_galerkin : sx_cell_galerkin;
const auto& sx_ey = (ey_type[0] == NODE) ? sx_node : sx_cell;
const int j_ex = (ex_type[0] == NODE) ? j_node_v : j_cell_v;
const int j_ey = (ey_type[0] == NODE) ? j_node : j_cell;
\end{codeblock}

这段源码有三个关键点。

\begin{enumerate}
\item \code{x = (xp - xyzmin.x)*dinv.x} 把物理坐标变成网格坐标。后面形函数都在无量纲网格坐标上计算。
\item \cpath{x} 和 \code{x - 0.5_rt} 分别对应 node-centered 和 cell-centered 自由度。也就是说，场分量的 staggered center 不是后处理标签，而是直接改变粒子看到的插值权重。
\item Galerkin 路径给 \cpath{ex/by/bz} 使用 \cpath{compute\_shape\_factor\_galerkin}，阶数是 \code{depos_order - 1}；非 Galerkin 时第二个模板参数为 0，因此阶数不变。
\end{enumerate}

以 2D XZ 编译为例，真正累加网格场的源码在 \cpath{Source/Particles/Gather/FieldGather.H}：

% <-- fenced code (cpp)
\begin{codeblock}
#elif defined(WARPX_DIM_XZ)
    // Gather field on particle Eyp from field on grid ey_arr
    for (int iz=0; iz<=depos_order; iz++){
        for (int ix=0; ix<=depos_order; ix++){
            Eyp += sx_ey[ix]*sz_ey[iz]*
                ey_arr(lo.x+j_ey+ix, lo.y+l_ey+iz, 0, 0);
        }
    }
    // Gather field on particle Exp from field on grid ex_arr
    // Gather field on particle Bzp from field on grid bz_arr
    for (int iz=0; iz<=depos_order; iz++){
        for (int ix=0; ix<=depos_order-galerkin_interpolation; ix++){
            Exp += sx_ex[ix]*sz_ex[iz]*
                ex_arr(lo.x+j_ex+ix, lo.y+l_ex+iz, 0, 0);
            Bzp += sx_bz[ix]*sz_bz[iz]*
                bz_arr(lo.x+j_bz+ix, lo.y+l_bz+iz, 0, 0);
        }
    }
    // Gather field on particle Ezp from field on grid ez_arr
    // Gather field on particle Bxp from field on grid bx_arr
    for (int iz=0; iz<=depos_order-galerkin_interpolation; iz++){
        for (int ix=0; ix<=depos_order; ix++){
            Ezp += sx_ez[ix]*sz_ez[iz]*
                ez_arr(lo.x+j_ez+ix, lo.y+l_ez+iz, 0, 0);
            Bxp += sx_bx[ix]*sz_bx[iz]*
                bx_arr(lo.x+j_bx+ix, lo.y+l_bx+iz, 0, 0);
        }
    }
    // Gather field on particle Byp from field on grid by_arr
    for (int iz=0; iz<=depos_order-galerkin_interpolation; iz++){
        for (int ix=0; ix<=depos_order-galerkin_interpolation; ix++){
            Byp += sx_by[ix]*sz_by[iz]*
                by_arr(lo.x+j_by+ix, lo.y+l_by+iz, 0, 0);
        }
    }
\end{codeblock}

公式上，第一段就是

\[
E_{y,p}
=
\sum_{i=0}^{p}\sum_{k=0}^{p}
S^{(x)}_{E_y,i} S^{(z)}_{E_y,k}
E_y(j_{E_y}+i,l_{E_y}+k).
\]

\cpath{Exp/Bzp}、\cpath{Ezp/Bxp}、\cpath{Byp} 的循环上限不同，是因为 Galerkin 插值会沿某些方向把 shape order 从 $p$ 降到 $p-1$。这不是任意优化，而是和离散 Maxwell operator、field staggering 与能量/电荷性质匹配的插值选择。

RZ 编译下，gather 先得到柱坐标分量，再转回笛卡尔粒子 pusher 需要的 $E_x,E_y,B_x,B_y$。关键转换在 \cpath{Source/Particles/Gather/FieldGather.H}：

% <-- fenced code (cpp)
\begin{codeblock}
    amrex::Real costheta;
    amrex::Real sintheta;
    if (rp > 0.) {
        costheta = xp/rp;
        sintheta = yp/rp;
    } else {
        costheta = 1.;
        sintheta = 0.;
    }
    const Complex xy0 = Complex{costheta, -sintheta};
    Complex xy = xy0;

    for (int imode=1 ; imode < n_rz_azimuthal_modes ; imode++) {

        // Gather field on particle Ethetap from field on grid ey_arr
        for (int iz=0; iz<=depos_order; iz++){
            for (int ix=0; ix<=depos_order; ix++){
                const amrex::Real dEy = (+ ey_arr(lo.x+j_ey+ix, lo.y+l_ey+iz, 0, 2*imode-1)*xy.real()
                                         - ey_arr(lo.x+j_ey+ix, lo.y+l_ey+iz, 0, 2*imode)*xy.imag());
                Ethetap += sx_ey[ix]*sz_ey[iz]*dEy;
            }
        }
        // Gather field on particle Erp from field on grid ex_arr
        // Gather field on particle Bzp from field on grid bz_arr
        for (int iz=0; iz<=depos_order; iz++){
            for (int ix=0; ix<=depos_order-galerkin_interpolation; ix++){
                const amrex::Real dEx = (+ ex_arr(lo.x+j_ex+ix, lo.y+l_ex+iz, 0, 2*imode-1)*xy.real()
                                         - ex_arr(lo.x+j_ex+ix, lo.y+l_ex+iz, 0, 2*imode)*xy.imag());
                Erp += sx_ex[ix]*sz_ex[iz]*dEx;
                const amrex::Real dBz = (+ bz_arr(lo.x+j_bz+ix, lo.y+l_bz+iz, 0, 2*imode-1)*xy.real()
                                         - bz_arr(lo.x+j_bz+ix, lo.y+l_bz+iz, 0, 2*imode)*xy.imag());
                Bzp += sx_bz[ix]*sz_bz[iz]*dBz;
            }
        }
        // Gather field on particle Ezp from field on grid ez_arr
        // Gather field on particle Brp from field on grid bx_arr
        for (int iz=0; iz<=depos_order-galerkin_interpolation; iz++){
            for (int ix=0; ix<=depos_order; ix++){
                const amrex::Real dEz = (+ ez_arr(lo.x+j_ez+ix, lo.y+l_ez+iz, 0, 2*imode-1)*xy.real()
                                         - ez_arr(lo.x+j_ez+ix, lo.y+l_ez+iz, 0, 2*imode)*xy.imag());
                Ezp += sx_ez[ix]*sz_ez[iz]*dEz;
                const amrex::Real dBx = (+ bx_arr(lo.x+j_bx+ix, lo.y+l_bx+iz, 0, 2*imode-1)*xy.real()
                                         - bx_arr(lo.x+j_bx+ix, lo.y+l_bx+iz, 0, 2*imode)*xy.imag());
                Brp += sx_bx[ix]*sz_bx[iz]*dBx;
            }
        }
        // Gather field on particle Bthetap from field on grid by_arr
        for (int iz=0; iz<=depos_order-galerkin_interpolation; iz++){
            for (int ix=0; ix<=depos_order-galerkin_interpolation; ix++){
                const amrex::Real dBy = (+ by_arr(lo.x+j_by+ix, lo.y+l_by+iz, 0, 2*imode-1)*xy.real()
                                         - by_arr(lo.x+j_by+ix, lo.y+l_by+iz, 0, 2*imode)*xy.imag());
                Bthetap += sx_by[ix]*sz_by[iz]*dBy;
            }
        }
        xy = xy*xy0;
    }

    // Convert Erp and Ethetap to Ex and Ey
    Exp += costheta*Erp - sintheta*Ethetap;
    Eyp += costheta*Ethetap + sintheta*Erp;
    Bxp += costheta*Brp - sintheta*Bthetap;
    Byp += costheta*Bthetap + sintheta*Brp;
\end{codeblock}

所以 \cpath{PushPX()} 里的 pusher 始终看见 Cartesian-like 的 \code{Exp,Eyp,Ezp,Bxp,Byp,Bzp}。RZ 的 Fourier mode 和柱坐标细节被 gather 层封装，但不是消失：它们决定粒子场的实际插值值。

把这层再向下看，WarpX 的 gather 还不是“只剩 3D/XZ/RZ 三个分支”。\cpath{FieldGather.H} 后面还继续给出了：

\begin{itemize}
\item \cpath{1D\_Z}：只沿 \cpath{z} 一个方向做 shape 累加，但仍保留 \cpath{Ez/Bx/By} 这组可能走 \cpath{p-1} 的 Galerkin 降阶路径；
\item \cpath{RCYLINDER}：先 gather \cpath{Fr/Ftheta/Fz}，再按粒子极角转回 \cpath{Fx/Fy/Fz}；
\item \cpath{RSPHERE}：先 gather \cpath{Er/Etheta/Ephi}，再按球坐标角转回 \cpath{Ex/Ey/Ez}；
\item \cpath{3D}：完整保留 \cpath{Ex/Ey/Ez/Bx/By/Bz} 六个分量各自不同的循环上限，因此 Galerkin 不是一个全局“降一阶开关”，而是对特定分量沿特定方向降阶。
\end{itemize}

这也能更准确地解释官方文档里 \cpath{energy-conserving} 与 \cpath{momentum-conserving} gather 的差别。\cpath{doGatherShapeN()} 本身并不读一个“gather family”枚举，它真正消费的是传进来的 \cpath{IndexType}。因此两族 gather 的代码差异不在这个 wrapper 里，而在更早的 field centering 上：

\begin{itemize}
\item \cpath{energy-conserving}：直接从原来的 staggered 或 nodal 网格 gather；
\item \cpath{momentum-conserving}：先把场中心化到 node，再把 nodal \cpath{IndexType} 送进 \cpath{doGatherShapeN()}。
\end{itemize}

这也解释了为什么 collocated grid 下两者等价，而 hybrid grid 下必须强制 \cpath{momentum-conserving}。一旦所有分量都已经是 nodal/collocated，\cpath{doGatherShapeN()} 看到的就只是同一组 \cpath{NODE/CELL} 组合。

如果只停在这层实现差异，这个问题还是讲浅了。Birdsall-Langdon 第一分卷 Chapter 10 给出的更硬边界是：\cpath{energy-conserving} 和 \cpath{momentum-conserving} 从来都不是同一套离散系统上“换一种 gather 插值”这么简单，而是两套不同的守恒合同。\cpath{momentum-conserving} 路线保留的是更自然的零总力/粒子-粒子相互作用对称结构，但它一般并不存在严格守恒的总能量；\cpath{energy-conserving} 路线则把

\[
W_E=\frac{V_c}{2}\sum_j \rho_j\phi_j
\]

当成第一性对象，再把粒子受力理解成

\[
\mathbf F_i=-\frac{\partial W_E}{\partial \mathbf x_i},
\]

也就是不再先“在网格上差分 $\phi$ 得 \cpath{E}、再把 \cpath{E} gather 给粒子”，而是要求粒子受力、离散 Poisson 解和场能量账本共享同一套 reciprocity 合同。对 WarpX 来说，这当然不意味着当前 \cpath{FieldGather.H} 里就直接复现了 Birdsall 的整个 energy-conserving electrostatic 变分算法；更准确的说法是，官方文档和 regression 里保留下来的 \code{energy-conserving gather} / \code{momentum-conserving gather} 命名，只有放回这条更老的理论分叉里才不会被误解成“两个 wrapper 里哪一个 stencil 更光滑”。

因此，本章后面凡是提到 gather family、field centering、collocated grid、Langmuir 守恒基线时，都应该记住自己实际上在比较三层东西：

\begin{enumerate}
\item \cpath{IndexType} 与 stagger/nodal centering 的实现差别；
\item sampled field 怎样被回插到粒子；
\item 这套回插究竟服务于哪一种离散守恒合同。
\end{enumerate}

还有一层容易漏掉：implicit path 并不复用显式 gather。\cpath{FieldGather.H} 的 \cpath{doGatherShapeNImplicit(...)} 会先按沉积算法分派：

\begin{itemize}
\item \cpath{Esirkepov}：\cpath{doGatherShapeNEsirkepovStencilImplicit}
\item \cpath{Villasenor}：\cpath{doGatherPicnicShapeN}
\item \cpath{Direct}：才退回普通半步位置 gather
\end{itemize}

所以 implicit 里不是“先统一 gather，再换 deposition”，而是 gather stencil 本身就要和后面的 charge-conserving 轨道几何解释匹配。

最后，external particle fields 也不止一种接入方式。\cpath{PushPX()} 里真实顺序是：

\begin{enumerate}
\item 先把常量 \cpath{m\_E\_external\_particle/m\_B\_external\_particle} 加到粒子局部 \cpath{Exp...Bzp}
\item 再对主网格场调用 \cpath{doGatherShapeN()}
\item 最后再执行 \cpath{GetExternalEBField} functor
\end{enumerate}

其中：

\begin{itemize}
\item \cpath{parse\_*\_ext\_particle\_function} 和 \cpath{repeated\_plasma\_lens} 是逐粒子直接相加；
\item \cpath{read\_from\_file} 则不会在 gather kernel 内逐粒子加场，而是先把外场加到 \cpath{Efield\_aux/Bfield\_aux} 这类 \cpath{MultiFab}，再让粒子像 gather 主场一样去 gather 它们。
\end{itemize}

对应的 regression 也已经有了比较清楚的分层：

\begin{itemize}
\item \cpath{energy\_conserving\_thermal\_plasma}：在 electrostatic 周期热等离子体里检查总能量增长不超过 \codeesc{0.3\%}，这是当前对 \code{energy-conserving gather} 最直接的物理断言；
\item \cpath{langmuir\_multi\_psatd\_momentum\_conserving}：在 \code{PSATD + momentum-conserving gather} 组合下继续用 Langmuir 问题做守恒/稳定性基线；
\item \cpath{load\_external\_field}：分别用 3D/RZ 磁镜单粒子轨道和时间缩放脚本验证 external particle fields 的 read-from-file、time dependency 和 multi-field superposition。
\end{itemize}

\section{位置推进与无质量粒子}

显式位置推进的实现位于 \cpath{Source/Particles/Pusher/UpdatePosition.H} 的 \cpath{GetExplicitPusherDisplacement()} 与 \cpath{UpdatePosition()}。\cpath{PhysicalParticleContainer::Evolve()} 的顺序是先调用 \cpath{doParticleMomentumPush(...)}，再在 \cpath{PositionPushType::Full} 分支调用 \cpath{UpdatePosition(...)}；因此这里消费的是推进后的时间中心动量，而不是另一个独立导出的速度数组。对有质量粒子，源码先计算

\[
\gamma^{-1}=\frac{1}{\sqrt{1+|\mathbf{u}|^2/c^2}},
\]

再按编译维度更新坐标：

\[
x\leftarrow x+u_x\gamma^{-1}\Delta t,
\qquad
y\leftarrow y+u_y\gamma^{-1}\Delta t,
\qquad
z\leftarrow z+u_z\gamma^{-1}\Delta t.
\]

对无质量粒子，源码使用

\[
\mathbf{v}=c\frac{\mathbf{u}}{|\mathbf{u}|}
\]

更新位置。因此 photon container 可以复用位置推进形式，但动量和沉积行为不同；光子容器的专门逻辑后续多物理章节再展开。

这个调用顺序也限定了如何验证“半步速度”。\cpath{UpdatePosition.H} 的注释把显式位置更新写成 $x(t+\Delta t)=x(t)+v(t+\Delta t/2)\Delta t$，但常规 Full 输出通常稳定提供的是位置和机械动量，而不是独立的 half-step velocity attribute。因此，相邻输出位置差可以构造与时间中心速度比较的 proxy，却不能被说成直接测得半步速度。对读者而言，正确的验证问题是：你比较的动量究竟处在哪个时间层，还是仅用相邻位置构造了平均量？

另一个容易忽略的分叉在 \cpath{PushSelector.H}：Boris 接受 \cpath{FirstHalf/SecondHalf/Full} 的 \cpath{momentum\_push\_type}，而当前 \cpath{UpdateMomentumHigueraCary()} 接口没有这一参数。因此不能把两个 pusher 写成完全相同的 split-half 输出合同。做轨道或时间层诊断时，先在输入中确认 pusher 和 push type，再比较位置、动量与场的时间层；只有这三者对齐，误差才可被归因于算法，而不是采样时刻不同。

\section{RR、implicit 与 photon path}

前面 4.1 到 4.10 主要还是显式带电粒子的主线，但 WarpX 的粒子推进并不只有这一条路径。至少还有三条不能忽略的分支：

\begin{enumerate}
\item classical radiation reaction；
\item implicit particle push；
\item photon container 的无质量推进。
\end{enumerate}

\textbf{读者的多物理状态交接卡。} 进入这一节以及后续的 ionization、collisions 和 QED 前，先区分“事件何时提交”与“新状态何时被 solver 消费”。外层时间步的固定入口顺序是：\cpath{doFieldIonization()} -> QED event pass -> \cpath{particleinjection} -> \cpath{OneStep()}。因此，已经在 field ionization 或 QED event pass 中创建的带电 product 会进入随后的 \cpath{OneStep()}；但这条顺序不能被误读为所有多物理过程都在 outer loop 的同一个位置发生。

\begin{enumerate}
\item \textbf{ionization。} ADK event 同时提高 source ion 的 \cpath{ionizationLevel} 并创建 electron product；随后由具体 solver 分支推进，并在 charged-particle deposition 中把离化态转成有效 source charge。应同时查看离子离化态、电子 product 以及 $\rho,\mathbf J$，不能只数新电子。
\item \textbf{QED。} \cpath{doQEDEvents()} 只消费已经由 pusher 演化到触发条件的 optical depth；它不是在 event pass 内从零开始抽样。Quantum Synchrotron 会保留并改写 lepton，Breit-Wheeler 会使 photon source 失效并创建 charged pairs；两者的 product 虽都在随后 \cpath{OneStep()} 前出现，却有不同的 source 命运。
\item \textbf{collisions。} \cpath{collisions.split\_momentum\_push} 只在 explicit 路径中组织半步动量：第一半步不沉积，collision 改写中间状态，第二半步和完整位置推进才交给正常沉积。关闭该选项时，collision 在完整 particle push 前执行；\cpath{OneStep\_JRhom()} 与 \cpath{OneStep\_sub1()} 都要求关闭 split momentum push。它是求解器内的时间调度，不是 outer-loop event pass。
\item \textbf{implicit。} 一次 nonlinear trial 不是新的物理外层步。\cpath{x\_n/ux\_n} 是 step-start reference state，收敛后的轨道及其 source 才属于该步的可解释结果；suborbit fallback 还会改变部分粒子的沉积路线。
\item \textbf{photon。} photon container 能演化位置和 Breit-Wheeler optical depth，却自身不沉积 charge/current。只有 photon 被转换后出现的 charged product 才可通过随后的 charged-particle 路径进入 $\rho/\mathbf J$。
\end{enumerate}

这张卡给出本节的验证尺度：RR 看解析动量/能量损失；ionization 看离化态、product 和有效 source；collision 看动量、product 与守恒或平衡量；QED 还要把 source 命运、optical depth 和 product 守恒一起比较。任何单一粒子数或单条轨道都不足以替代这些成组观察量。

先看 RR。\cpath{Source/Particles/Pusher/PushSelector.H} 的 \cpath{doParticleMomentumPush()} 说明，RR 不是第四种独立 pusher，而是优先级高于 \cpath{ParticlePusherAlgo} 的一个分支。省略编译宏后的控制流可概括为：

% <-- fenced code (cpp)
\begin{codeblock}
if (do_crr) {
    if (qed_sync && chi >= t_chi_max) {
        UpdateMomentumBoris(...);
    } else {
        UpdateMomentumBorisWithRadiationReaction(...);
    }
} else if (pusher_algo == ParticlePusherAlgo::Boris) {
    UpdateMomentumBoris(...);
} else if (pusher_algo == ParticlePusherAlgo::Vay) {
    UpdateMomentumVay(...);
} else if (pusher_algo == ParticlePusherAlgo::HigueraCary) {
    UpdateMomentumHigueraCary(...);
}
\end{codeblock}

因此，打开 \cpath{do\_crr} 后，当前粒子不会再走 Vay 或 Higuera--Cary，而是进入 Boris 家族分支。通常它调用“Boris 加辐射反作用”；但在编译 QED 且同步开关生效时，代码会先计算 $\chi$：$\chi<t_{\chi,\max}$ 才调用该 RR 例程，较高 $\chi$ 则调用普通 Boris。这一门限的含义是：源码保证了 pusher 家族的选择，却不保证每一个 $\chi$ 都附加 classical RR。\cpath{Source/Particles/Pusher/UpdateMomentumBorisWithRadiationReaction.H} 中的 \cpath{UpdateMomentumBorisWithRadiationReaction()} 则表明低 $\chi$ 分支如何实现：它先调用普通 \cpath{UpdateMomentumBoris()}，再用新旧动量平均构造中间时刻的 $\gamma_n$、$\mathbf{v}_n$ 和 Lorentz force，最后再把辐射反作用力乘 \cpath{dt} 加回动量。代码结构上，这是一种 Boris 后附加阻尼项，而不是完全重写一套 relativistic mover。

再看 implicit path。它和显式 \cpath{PushPX()} 的根本区别，不在于换了另一个 \cpath{UpdateMomentum*()}，而在于时间层和收敛逻辑都改了。\cpath{Source/Particles/Pusher/ImplicitPushPX.cpp} 的注释直接说明了顺序：

\begin{enumerate}
\item 先 position push 半步；
\item 再 gather 场；
\item 再做 velocity push；
\item 再把 old/new velocity 平均成 time-centered 值；
\item 位置和速度彼此依赖，因此做 Picard 固定点迭代，直到 step norm 收敛。
\end{enumerate}

而这里真正把上一篇属性图接进来的，是 \cpath{x\_n/y\_n/z\_n}、\cpath{ux\_n/uy\_n/uz\_n} 和 \cpath{nsuborbits}。在 \cpath{Source/Particles/Pusher/ImplicitPushPX.cpp}，这些量被明确当成“the positions and velocities saved at the start of the step”取出；随后粒子初值直接从 \cpath{x\_n} 和 \cpath{ux\_n} 开始，而不是从当前位置盲目继续推进。也就是说，\cpath{x\_n/ux\_n} 在 implicit 路径里不是诊断缓存，而是 nonlinear solve 的参考态。

\cpath{nsuborbits} 则是 implicit 不收敛时的 fallback 状态。\cpath{ImplicitPushXP()} 在 \cpath{Source/Particles/Pusher/ImplicitPushPX.cpp} 中，如果粒子没收敛，就把 \code{nsuborbits[ip] = 2}，再通过 \cpath{SetupSuborbitParticles()} 把这些粒子的权重临时置零、单独收集索引。后续 \cpath{ImplicitPushXPSubOrbits()} 又会强制把沉积算法切到 Villasenor，见 \cpath{ImplicitPushPX.cpp}。所以 suborbit 不只是“多分几步时间步”，还会连带改变当前粒子的沉积路径。

最后看 photon container。\cpath{Source/Particles/PhotonParticleContainer.cpp} 的 \cpath{PhotonParticleContainer::Evolve()} 并没有重写 species 外层循环，而是继续调用 \cpath{PhysicalParticleContainer::Evolve(...)}。也就是说，tile loop、AMR buffer 分区、gather 外壳这些基础设施仍然复用。但 photon 通过两层专门改写改变了物理语义：

\begin{itemize}
\item \cpath{PhotonParticleContainer.H} 中 \cpath{DepositCharge()} 和 \cpath{DepositCurrent()} 都是空实现；
\item \cpath{PhotonParticleContainer::PushPX()} 在 \cpath{Source/Particles/PhotonParticleContainer.cpp} 里只做 gather、可选 Breit-Wheeler optical depth 演化、以及无质量 \cpath{UpdatePosition(...)}，并不调用 Boris/Vay/Higuera-Cary 的带电动量更新。
\end{itemize}

因此 photon path 和普通 charged species 的差异不只是“不沉积电流”，而是连 momentum update 的物理模型都变了。它消费的 runtime attributes 主要是：

\begin{itemize}
\item builtin \cpath{ux/uy/uz}；
\item \cpath{opticalDepthBW}；
\item 可选 \cpath{*\_btd}；
\end{itemize}

而不会消费普通 charged implicit 路径里的 \cpath{ionizationLevel}、\cpath{opticalDepthQSR}、\cpath{x\_n/ux\_n/nsuborbits}。

从这一层往后看，WarpX 粒子推进可以总结成三种结构：

\begin{itemize}
\item 显式主线：Boris / Vay / Higuera-Cary；
\item 显式修正：Boris + classical RR；
\item 非标准推进：implicit fixed-point / suborbit fallback / photon push。
\end{itemize}

它们共享的是 \cpath{PhysicalParticleContainer::Evolve()}、gather 外壳和粒子属性系统；真正分叉的是时间层、收敛控制、沉积算法约束和具体消费的 runtime attributes。

如果再往 implicit solver 深处走一步，还要继续把“suborbit 轨道本身”和“JFNK 线性化源项拼装”分开看。\cpath{Source/FieldSolver/ImplicitSolvers/ImplicitSolver.cpp} 明确把 linear stage 的电流写成

\[
J(E)=J_{\mathrm{suborbit}}+J_0+\mathrm{MM}(E-E_0).
\]

这里：

\begin{itemize}
\item \cpath{J\_0} 对应 \cpath{current\_fp\_non\_suborbit}；
\item \cpath{MM} 对应 \cpath{MassMatrices\_X/Y/Z}；
\item \cpath{J\_suborbit} 才是那些真正需要 suborbit fallback 的粒子继续显式推进后沉到 \cpath{current\_fp} 的部分。
\end{itemize}

这正好解释了为什么 \cpath{MultiParticleContainer::Evolve()} 在 implicit 模式下还要额外处理 \cpath{current\_fp\_non\_suborbit} 和 \cpath{MassMatrices\_PC} 的清零时机，见 \cpath{Source/Particles/MultiParticleContainer.cpp}。它不是普通 bookkeeping，而是在维护 JFNK 的三项分拆。

WarpX 为这一节提供了一条直接的 regression 入口：\cpath{Examples/Tests/radiation\_reaction/}。它不是应用级 checksum，而是强 analysis：

\begin{itemize}
\item 平行动量 case 要求 \cpath{gamma} 保持不变；
\item 垂直动量 case 要求 \cpath{gamma(t)} 满足解析 Landau-Lifshitz 衰减公式；
\item 容差统一为 \codeesc{5\%}
\end{itemize}

因此它正好锚定了上面这条“先 Boris，再加 RR 修正”的源码路径，而不是泛化地证明“高能粒子大概会辐射”。

\cpath{ImplicitPushXPSubOrbits()} 里还有两条实现约束非常重要。第一，\cpath{Source/Particles/Pusher/ImplicitPushPX.cpp} 强制把 suborbit 路径的 current deposition 切到 Villasenor：

% <-- fenced code (cpp)
\begin{codeblock}
const auto depos_type = CurrentDepositionAlgo::Villasenor;
\end{codeblock}

所以一旦粒子进入 suborbit fallback，用户原来选择的沉积算法并不会继续沿用。第二，\cpath{Source/Particles/Pusher/ImplicitPushPX.cpp} 把 \cpath{deposit\_mass\_matrices} 限定成

% <-- fenced code (cpp)
\begin{codeblock}
use_mass_matrices_pc && !linear_stage_of_jfnk
\end{codeblock}

这说明 suborbit 粒子的 mass-matrix deposit 当前只服务 preconditioner，不直接服务 Jacobian 的 linear stage。

于是 implicit 粒子推进实际上还要再分成四层：

\begin{enumerate}
\item \cpath{x\_n/ux\_n/nsuborbits} 这组 step-start 属性；
\item \cpath{PushXPSingleStep()} 的 fixed-point 单轨道求解；
\item 不收敛粒子的 suborbit fallback 与 Villasenor-only 重沉积；
\item \code{current_fp_non_suborbit + MM*(E-E_0) + J_suborbit} 这套线性化源项拼装。
\end{enumerate}

这也是为什么 implicit 路径不能被简单概括成“把显式 pusher 改成 Picard 迭代”。它已经把粒子推进、沉积算法、质量矩阵和 Newton/JFNK 线性化绑成了一套更大的数值系统。

\section{Field Ionization：ADK 不是附加后处理，而是粒子属性和沉积权重一起改写}

前面已经出现过 \cpath{ionizationLevel}，但如果不单独拆开，很容易把 field ionization 理解成“额外生成几个电子”。源码里真实发生的是三件事一起改变：

\begin{enumerate}
\item 源离子 species 在运行时新增整数属性 \cpath{ionizationLevel}；
\item 电离事件通过 \cpath{filterCopyTransformParticles} 复制到 product electron species；
\item 后续 \cpath{rho/J} 沉积时，源离子的有效电荷按 \cpath{ionizationLevel} 动态乘上去。
\end{enumerate}

\cpath{PhysicalParticleContainer::InitIonizationModule()} 会在拿到运行时 \cpath{dt} 后，才真正初始化 ADK 所需的：

\begin{itemize}
\item \cpath{ionization\_energies}
\item \cpath{adk\_power}
\item \cpath{adk\_prefactor}
\item \cpath{adk\_exp\_prefactor}
\item 可选 \cpath{adk\_correction\_factors}
\end{itemize}

同时还会在没有现成属性时补上：

% <-- fenced code (cpp)
\begin{codeblock}
if (!HasiAttrib("ionizationLevel")) {
    AddIntComp("ionizationLevel");
}
\end{codeblock}

这说明 \cpath{ionizationLevel} 属于“模块初始化阶段动态加入的持久物理状态”，而不是 \cpath{PIdx} builtin。

这里还有两个容易被漏掉的源码边界。

第一，WarpX 的 field ionization 只有 \cpath{ADK} 主链，没有独立 \cpath{OTB} 分支。官方参数文档对 \cpath{do\_field\_ionization} 的描述明确写的是 \code{using the ADK theory}，源码里读到的也只有 \cpath{do\_adk\_correction}、\cpath{adk\_prefactor}、\cpath{adk\_exp\_prefactor} 和 \cpath{adk\_power}。因此如果后面讨论 \cpath{OTB}，那只能作为“源码未实现的外部模型边界”，不能写成 \cpath{Ionization.*} 已有的一条并行实现。

第二，\cpath{physical\_element} 也不是“用户手工给一串电离能表”的接口。\cpath{InitIonizationModule()} 实际是通过 \cpath{ion\_map\_ids}、\cpath{ion\_atomic\_numbers} 和 \cpath{ion\_energy\_offsets} 从 WarpX 内建的 \cpath{table\_ionization\_energies} 里切出整段 successive ionization energies，再按当前运行时 \cpath{dt} 现算 ADK prefactor。因此 species 构造期不能完成这一步，真正原因不是代码组织习惯，而是 ADK 率已经把本步 \cpath{dt} 吸进了系数里。

真正的单粒子电离判定在 \cpath{Particles/ElementaryProcess/Ionization.H} 的 \cpath{IonizationFilterFunc::operator()}。它不会只看实验室系 \code{|E|}，而是先 gather 主场和 external particle field，再结合当前 \cpath{ux/uy/uz} 计算粒子系场幅值，然后才用 ADK 公式给出概率：

% <-- fenced code (cpp)
\begin{codeblock}
Real w_dtau = (E <= 0._rt) ? 0._rt : 1._rt/ ga * m_adk_prefactor[ion_lev] *
    std::pow(E, m_adk_power[ion_lev]) *
    std::exp( m_adk_exp_prefactor[ion_lev]/E );
const Real p = 1._rt - std::exp( - w_dtau );
\end{codeblock}

因此 boosted-frame \cpath{field\_ionization} regression 不是“换个坐标跑一遍”而已，它实际上在验证这层相对论一致性。

更容易被忽略的是 transform 本身。\cpath{IonizationTransformFunc} 看起来只有一句：

% <-- fenced code (cpp)
\begin{codeblock}
src.m_runtime_idata[0][i_src] += 1;
\end{codeblock}

但这并不意味着 product electron 没被创建。真正的 product species 写入，是由 \cpath{MultiParticleContainer::doFieldIonization()} 里的：

% <-- fenced code (cpp)
\begin{codeblock}
filterCopyTransformParticles<1>(...)
\end{codeblock}

统一完成的。也就是说，一次电离事件被拆成：

\begin{itemize}
\item 源离子 \code{ionizationLevel += 1}
\item 新电子复制到 \cpath{ionization\_product\_species}
\end{itemize}

两部分。

最后，这个离化态不会只停留在 diagnostics。\cpath{WarpXParticleContainer::DepositCurrent()} 和 \cpath{DepositCharge()} 都会把 \cpath{ionizationLevel} 传下去，而底层沉积 kernel 会做：

% <-- fenced code (cpp)
\begin{codeblock}
if (do_ionization) { wq *= ion_lev[ip]; }
\end{codeblock}

因此 field ionization 的真正闭环是：

\[
\text{ADK event}
\rightarrow \texttt{ionizationLevel}
\rightarrow \text{effective source charge}
\rightarrow \rho,J
\rightarrow \text{fields and diagnostics}.
\]

从粒子推进的视角看，这说明 WarpX 的多物理不总是“在主循环旁边加一个模块”。像 field ionization 这样的过程，会直接改写粒子属性系统和沉积权重，所以它应该被视为 particle algorithm 本体的一部分。

更细一点说，这里的“改写沉积权重”并不是去改宏粒子统计权重 \cpath{w}，而是保持 \cpath{w} 不变、只把有效物理电荷改成

\[
w q_e \times \texttt{ionizationLevel}.
\]

这也是 \cpath{InitIonizationModule()} 一开始就把 ionizable species 的 \cpath{charge} 强制改回 \cpath{q\_e} 的原因：源码要把“多少个 physical particles”与“每个 physical particle 当前带几个基本电荷”这两层语义拆开，前者留在 \cpath{w}，后者留在 \cpath{ionizationLevel}。

\section{Collisions：不是旁路模块，而是和 momentum push 强耦合的时间调度层}

field ionization 还是“每个源 species 自己带 product species 逻辑”的模块，但 collision 从入口层开始就不是这样。\cpath{MultiParticleContainer} 构造时直接建：

% <-- fenced code (cpp)
\begin{codeblock}
collisionhandler = std::make_unique<CollisionHandler>(this);
\end{codeblock}

运行时则统一走：

% <-- fenced code (cpp)
\begin{codeblock}
collisionhandler->doCollisions(step, cur_time, dt, this);
\end{codeblock}

这说明 collision 在 WarpX 里的第一抽象层不是某个 species 的附加属性，而是一个面向整个 \cpath{MultiParticleContainer} 的统一调度器。

\cpath{CollisionHandler} 的第一职责也不是实现物理公式，而是把：

% <-- fenced code (text)
\begin{codeblock}
collisions.collision_names
<collision_name>.type
\end{codeblock}

映射成具体对象。当前入口层已经能分出：

\begin{itemize}
\item \cpath{pairwisecoulomb}
\item \cpath{background\_mcc}
\item \cpath{pulsed\_decay}
\item \cpath{background\_stopping}
\item \cpath{dsmc}
\item \cpath{nuclearfusion}
\item \cpath{bremsstrahlung}
\item \cpath{linear\_breit\_wheeler}
\item \cpath{linear\_compton}
\end{itemize}

其中有些是独立的 \cpath{CollisionBase} 派生类，有些则走模板化的 \code{BinaryCollision<CollisionFunc, CreationFunc>}。因此“collision family”在 WarpX 里从一开始就包含了既可能只改动动量、也可能会生成新粒子的过程。

真正的公共合同在 \cpath{CollisionBase}。它提供的不是一条统一散射公式，而是：

\begin{itemize}
\item \cpath{species}
\item \cpath{ndt}
\item \codeesc{CollisionSteppingMode::\{Supercycle, Subcycle\}}
\end{itemize}

两种 stepping mode 的语义不同：

\begin{itemize}
\item \cpath{Subcycle}：一个 PIC step 内跑 \cpath{ndt} 次，每次用 \cpath{dt/ndt}
\item \cpath{Supercycle}：每 \cpath{ndt} 个 PIC steps 才跑一次，但单次碰撞用 \cpath{dt*ndt}
\end{itemize}

这说明 collision 入口首先定义的是“碰撞算子相对于 PIC 主步怎样调度”，而不是“碰撞算子本体长什么样”。

更关键的是，collision 从全局参数读取阶段就已经和 momentum pusher 强耦合。\cpath{WarpX.cpp} 只要看到：

% <-- fenced code (text)
\begin{codeblock}
collisions.collision_names = ...
\end{codeblock}

默认就会打开：

% <-- fenced code (cpp)
\begin{codeblock}
m_collisions_split_momentum_push = true;
\end{codeblock}

然后再允许用户用：

% <-- fenced code (text)
\begin{codeblock}
collisions.split_momentum_push
\end{codeblock}

覆盖。这说明 collision 默认并不是“在完整粒子推进前或后统一做一次”，而是插在 momentum push 的中间。源码还明确给了当前边界：

\begin{itemize}
\item implicit evolve scheme 下不真正支持 split momentum push
\item Higuera-Cary 目前也不支持
\end{itemize}

因此，collision 入口从一开始就已经和：

\begin{itemize}
\item pusher 选择
\item electrostatic / electromagnetic solver
\item 主循环时间组织
\end{itemize}

绑在一起。

这层耦合并不是只靠源码阅读猜出来的，\cpath{analysis\_test\_2d\_collisions\_split\_momentum\_push.py} 直接拿 reduced diagnostics 的：

\begin{itemize}
\item \cpath{field\_energy}
\item \cpath{particle\_energy}
\end{itemize}

做两类断言：

\begin{enumerate}
\item 总能量守恒误差必须足够小
\item 场能量涨落长期平均必须接近 equipartition 参考值
\end{enumerate}

所以这组 regression 真正在验证的是 collision insertion ordering 的数值合同，而不是某个散射截面是否长得对。

从粒子推进章节的视角看，这意味着 WarpX 的多物理插入至少分成三种类型：

\begin{enumerate}
\item 改写粒子属性和沉积权重：field ionization
\item 改写 momentum push 的时间组织：collisions
\item 改写单粒子推进与产品粒子统计：QED / photon emission / pair creation
\end{enumerate}

它们都不是“主循环旁边挂一个功能块”这么简单，而是直接进入粒子算法本体。

\subsection{\texttt{BinaryCollision} 先产出事件表，\texttt{ParticleCreationFunc} 再真正创建 products}

如果继续沿 collisions 这一支往下读，就会发现 \code{BinaryCollision<CollisionFunctor, CopyTransformFunctor>} 的 cell-level functor 并不直接往 product species 里塞粒子。它先输出的是几张事件表：

\begin{itemize}
\item \cpath{p\_mask}
\item \cpath{p\_pair\_indices\_1}
\item \cpath{p\_pair\_indices\_2}
\item \cpath{p\_pair\_reaction\_weight}
\end{itemize}

也就是：

\begin{itemize}
\item 哪一对 parent 真发生了反应
\item 这次反应对应哪两个 parent 粒子
\item 这次反应要抽走多少权重交给 products
\end{itemize}

真正把这些事件落到 product species tile 里的，是后面的 \cpath{ParticleCreationFunc}。

这层分工很重要，因为它说明 collision 模块内部也不是“一次 kernel 既判断又创建”，而是：

% <-- fenced code (text)
\begin{codeblock}
CollisionFunc
-> event tables
-> ParticleCreationFunc
-> type-specific momentum initialization
\end{codeblock}

\cpath{ParticleCreationFunc} 里最关键的分叉是：

% <-- fenced code (cpp)
\begin{codeblock}
const int products_per_reactant_factor =
    (m_collision_type == CollisionType::LinearCompton) ? 1 : 2;
\end{codeblock}

它把 binary creation 分成两类：

\begin{enumerate}
\item \cpath{LinearCompton}
\begin{itemize}
\item 每个 product species 每个事件只创建 1 个宏粒子
\item 散射 photon 继承入射 photon 位置
\item 散射 electron 继承入射 electron 位置
\end{itemize}
\item 几乎所有其他 binary creation
\begin{itemize}
\item 包括 \cpath{LinearBreitWheeler}
\item 都会在两个 reactant 的位置各复制一套 products
\end{itemize}
\end{enumerate}

这就是为什么 \cpath{LinearBreitWheeler} 虽然物理上只产生一对 \code{electron + positron}，实现上却会生成：

\begin{itemize}
\item 2 个电子宏粒子
\item 2 个正电子宏粒子
\end{itemize}

并且每个宏粒子只拿一半反应权重。WarpX 用这种“复制到两个 parent 位置”的实现换取了局域电荷守恒。

\cpath{LinearCompton} 则反过来是当前最明显的特例。它不会做双份复制，而是：

\begin{itemize}
\item 1 个散射 photon
\item 1 个散射 electron
\end{itemize}

然后用 \cpath{LinearComptonInitializeMomentum()} 在电子静止系里按 Klein-Nishina 分布抽样散射角，再 Lorentz 变换回 lab frame，并用总动量守恒补出散射后电子动量。

所以如果把 collisions 再往 product 创建层推进一步，当前最值得记的不是某个碰撞截面公式，而是：

\begin{itemize}
\item \cpath{BinaryCollision} 先做事件压缩
\item \cpath{ParticleCreationFunc} 再做统一 product 落地
\item \cpath{LinearBreitWheeler} 和 \cpath{LinearCompton} 在“一个事件到底创建几个宏粒子”这件事上故意采用了两种不同合同
\end{itemize}

这也解释了为什么 \cpath{linear\_breit\_wheeler} 与 \cpath{linear\_compton} 的 regression 都把：

\begin{itemize}
\item 守恒量
\item product species 数目
\item 最终产额
\end{itemize}

看得非常重，因为真正容易出错的地方就在 product 布局和权重分配，不只是事件概率。

\subsection{\texttt{BackgroundMCC}、\texttt{PulsedDecay} 与 DSMC：collision 还有三条完全不同的后半段}

前一小节已经把 \code{BinaryCollision -> ParticleCreationFunc} 这条 product 创建主链打通了，但 WarpX 的 collision 还至少有三条不能混写进去的实现分叉。

第一条是 \cpath{BackgroundMCCCollision}。它不是“两个显式 species 配对后再散射”，而是：

\begin{itemize}
\item 一个显式运动 species
\item 一个由 \code{background_density(x,y,z,t)}、\code{background_temperature(x,y,z,t)} 或常数描述的背景气体
\item 一个用 \cpath{m\_max\_background\_density} 和截面表扫出来的 \cpath{nu\_max}
\end{itemize}

因此它的核心上游稳定性门闩是

\[
\nu_{\max}\Delta t,
\]

而不是前面 \cpath{BinaryCollision} 那种 pair enumeration。更关键的是，若打开 impact ionization，它也不是走 \cpath{ParticleCreationFunc}，而是走：

% <-- fenced code (text)
\begin{codeblock}
ImpactIonizationFilterFunc
-> filterCopyTransformParticles
-> 原电子动量更新 + 新电子/新离子创建
\end{codeblock}

也就是说，\cpath{BackgroundMCC} 的 ionization 是 filter/copy/transform 分支，不是 pair-event-table 分支。

第二条是 \cpath{PulsedDecay}。它甚至不再是 pairwise 碰撞，而是 cell 内总权重衰变问题。源码先在每个 cell 统计 parent species 的总权重，再按

\[
W_{\text{prod}} = W_1\left(1-e^{-\nu(x,y,z,t)\Delta t}\right)
\]

算出本步应衰变掉多少权重，然后再用 \cpath{fixed\_product\_weight} 把这份总权重离散成 product 宏粒子数。新粒子的位置和基底速度来自该 cell 内随机挑选的 parent 粒子，再叠加方向相关 thermal speed。这里真正的物理合同是：

\begin{itemize}
\item parser 驱动的 \code{decay_rate(x,y,z,t)}
\item fixed-weight 宏粒子离散化
\item parent 权重逐步扣减
\end{itemize}

而不是“一对 parent 粒子 -> 一次散射事件”。

第三条是 DSMC 的 \cpath{SplitAndScatterFunc}。它虽然仍挂在 \cpath{BinaryCollision} 大框架下，但前半段 \cpath{DSMCFunc} 只负责：

\begin{itemize}
\item 组 pair
\item 调 \cpath{CollisionPairFilter}
\item 写 \cpath{p\_mask/p\_pair\_indices/p\_pair\_reaction\_weight}
\item 先从 reactants 扣权重
\end{itemize}

真正的后半段 product/child 创建则交给 \cpath{SplitAndScatterFunc}，而不是前面那条 \cpath{ParticleCreationFunc}。这里最重要的 slot 语义是：

\begin{itemize}
\item slot 0：第一 reactant 的 child copy
\item slot 1：第二 reactant 的 child copy
\item slot 2/3：true products
\end{itemize}

于是：

\begin{itemize}
\item elastic / back 只在 slot 0/1 生成 weight-split child，并在 COM frame 随机旋转或反转速度；
\item \cpath{charge\_exchange} / \cpath{two\_product\_reaction} 在 slot 2/3 生成 products，并由 \cpath{TwoProductComputeProductMomenta(...)} 统一处理动量；
\item \cpath{charge\_exchange} 还会交换 products 的生成位置，以维持局域电荷守恒；
\item DSMC ionization 则会同时保留 slot 0 的 incident child，并在 slot 2/3 生成 ejected electron 和 ion。
\end{itemize}

所以到这里，WarpX 的 collision 至少要分成四种后半段语义：

\begin{enumerate}
\item \cpath{BackgroundMCC} 的背景介质 + filter/transform
\item \cpath{PulsedDecay} 的 cell 内总权重衰变
\item \code{BinaryCollision + ParticleCreationFunc}
\item \code{BinaryCollision + SplitAndScatterFunc}
\end{enumerate}

这也是为什么不能再用一套统一口径把 \cpath{BackgroundMCC}、\cpath{PulsedDecay}、DSMC、\cpath{LinearCompton} 和 QED 都简单写成“会产生新粒子的碰撞”。

在 regression 层，这三条分叉也看完全不同的量：

\begin{itemize}
\item \cpath{analysis\_collision\_3d\_pulsed\_decay.py} 比的是最终 ion 总权重和 0D 衰变模型；
\item \cpath{analysis\_ionization\_dsmc\_3d.py} 比的是 \cpath{n\_e}、\cpath{n\_n} 和 \cpath{n\_eT\_e} 的全局模型；3D ion-impact sibling 则是 \code{ion_impact_ionization.dat + ions neutrals} 的 checksum-only 分叉。
\item \cpath{analysis\_charge\_exchange\_dsmc\_1d.py} 比的是通量指数衰减；
\item \cpath{analysis\_two\_product\_reaction\_dsmc\_1d.py} 比的是 products 的理论速度；
\item \cpath{analysis\_photoneutralization\_dsmc\_1d.py} 还把快电子的产物验证接到了 \cpath{BoundaryScraping} diagnostics。
\end{itemize}


\subsection{\texttt{pairwisecoulomb}、\texttt{bremsstrahlung}、\texttt{background\_stopping} 与 \texttt{nuclearfusion}}

把上一节的 \code{BackgroundMCC / PulsedDecay / DSMC} 再往外扩，WarpX 当前 collision 里还剩四条不能并进同一套 product 语义的主分叉。

第一条是 \cpath{pairwisecoulomb}。它虽然也走 \cpath{BinaryCollision} 的 cell 内 pairing 外壳，但 \cpath{PairWiseCoulombCollisionFunc} 最后只做一件事：

% <-- fenced code (cpp)
\begin{codeblock}
ElasticCollisionPerez(...)
\end{codeblock}

也就是说，它没有 products，也不依赖后面的 \cpath{ParticleCreationFunc} 或 \cpath{SplitAndScatterFunc}。构造期真正重要的只有：

\begin{itemize}
\item \cpath{CoulombLog}
\item \cpath{use\_global\_debye\_length}
\end{itemize}

当 \code{CoulombLog < 0} 且不使用 global Debye length 时，源码会进一步要求运行时 local temperature，从而按局域状态估算碰撞尺度。它的 regression 口径也因此不是“有没有新粒子”，而是：

\begin{itemize}
\item \cpath{analysis\_collision\_1d.py} 检查 relaxation benchmark
\item \cpath{analysis\_collision\_1d\_correct\_conservation.py} 检查总动量和总动能守恒
\end{itemize}

第二条是 \cpath{bremsstrahlung}。它重新回到 pair-event 表，但后半段 product 创建也不是通用 \cpath{ParticleCreationFunc}，而是专门的 \cpath{PhotonCreationFunc}。前半段 \cpath{BremsstrahlungFunc} 只给出：

\begin{itemize}
\item 哪一对发生事件
\item photon 权重
\item photon 能量
\end{itemize}

后半段再由 \cpath{PhotonCreationFunc}：

\begin{enumerate}
\item 在第一 reactant 位置上创建 photon
\item 同时回写 parent electron/ion 动量
\item 用能量动量守恒补全 photon 动量
\end{enumerate}

因此 \cpath{bremsstrahlung} 的真实语义是“先算失能和 photon energy，再做 parent+product 一起更新”。\cpath{analysis\_collision\_1d\_Bremsstrahlung.py} 也正是按这个口径验证：

\begin{itemize}
\item 总能量守恒
\item 总动量守恒
\item \cpath{dE/dx} 与解析估计
\item 每步新 photon 数与解析截面估计
\end{itemize}

第三条是 \cpath{background\_stopping}。这已经不再是二体散射，而是对单 species 应用解析 slowing-down law。源码在 \code{background_type = electrons} 与 \cpath{ions} 之间分成两条公式：

\begin{itemize}
\item 对电子背景：
\end{itemize}
\[
  \frac{d\mathbf{u}}{dt}=-\alpha\mathbf{u}
  \quad\Rightarrow\quad
  \mathbf{u}^{n+1}=\mathbf{u}^n e^{-\alpha\Delta t}
\]
\begin{itemize}
\item 对离子背景：
\end{itemize}
\[
  \frac{dW}{dt}=-\frac{\alpha}{\sqrt{W}}
\]
再按解析积分结果缩放速度

因此 \cpath{background\_stopping} 的 regression 口径也完全不同：\cpath{ion\_stopping/analysis.py} 直接用 Python 重写同一组 slowing-down 公式，对 constant / parsed 的电子与离子背景逐点对照最终粒子能量。

第四条是 \cpath{nuclearfusion}。它又回到 pair-event 框架，但事件概率控制比普通 \cpath{BinaryCollision} 更复杂。构造期最关键的不是单个截面文件，而是：

\begin{itemize}
\item \cpath{event\_multiplier}
\item \cpath{probability\_threshold}
\item \cpath{probability\_target\_value}
\item fusion type
\end{itemize}

其中 \cpath{event\_multiplier} 用来提高稀有 fusion 事件的统计量，而 \cpath{probability\_threshold} / \cpath{probability\_target\_value} 用来在单 pair 概率过大时自动把 multiplier 压回安全范围，避免系统性低估 yield。源码上的截面模型至少分成：

\begin{itemize}
\item \cpath{BoschHaleFusionCrossSection}：D-D / D-T 之类经典两产物 fusion
\item \cpath{ProtonBoronFusionCrossSection}：p-B11，对应 Tentori-Belloni 2023 的拟合
\end{itemize}

它们的 regression 也相应分层：

\begin{itemize}
\item \cpath{analysis\_two\_product\_fusion.py} 检查 reactant 消耗、product 数、能量动量守恒、各向同性和理论 yield
\item \cpath{analysis\_proton\_boron\_fusion.py} 进一步检查 p-B11 的三 alpha 产物、热率拟合，以及 \cpath{probability\_threshold} 过大时的故意失真场景
\end{itemize}

所以如果把已经成文的 collision 分支一起看，WarpX 至少已经出现了下面这些互不等价的后半段：

\begin{enumerate}
\item 纯动量散射：\cpath{pairwisecoulomb}
\item filter/transform：\cpath{BackgroundMCC} ionization
\item 固定权重衰变：\cpath{PulsedDecay}
\item DSMC child/product 重建：\cpath{SplitAndScatterFunc}
\item 通用 products：\cpath{ParticleCreationFunc}
\item photon 专用 products：\cpath{bremsstrahlung} 的 \cpath{PhotonCreationFunc}
\item fusion event-probability + specialized kinematics：\cpath{nuclearfusion}
\end{enumerate}

这说明到目前为止，“collision 模块”已经不能再被当作一个统一 product 创建器来讲，而必须按物理合同和后半段实现机制分开处理。

\subsection{kernel 细节层：Perez 有效碰撞密度、Bremsstrahlung photon 写回、fusion products 复制语义}

把上一节的四条分叉再往下读，最容易误判的不是类型分派，而是它们在 kernel 末端到底怎样把统计事件变成具体粒子动量。

对 \cpath{pairwisecoulomb} 来说，\cpath{PairWiseCoulombCollisionFunc} 真正交给 \cpath{ElasticCollisionPerez(...)} 的并不是“原始 cell 密度”，而是一套先闭合局域尺度、再构造有效配对密度的量：

\begin{itemize}
\item 若 \code{CoulombLog < 0}，先按局域 \code{n,T} 算 Debye 长度
\item 再取 \code{bmax = max(lambda_D, r_min)}，保证 screening length 不小于原子间距
\item 然后构造加权后的
\end{itemize}
\[
  n_{12}
\]
来驱动 Perez 散射

因此 WarpX 当前的 weighted macroparticle Coulomb 碰撞，真正进入单对散射 kernel 的强度不是简单的 cell 物理密度，而是“局域 plasma 尺度 + 抽样配对统计”共同决定的有效量。这也解释了为什么 Coulomb regression 里一个重点看 relaxation，另一个重点看守恒。

对 \cpath{bremsstrahlung} 来说，上一节只讲到了前半段 \cpath{BremsstrahlungFunc} 如何给出 photon energy。再往后看 \cpath{PhotonCreationFunc}，会发现它不是只把 energy 填进 product，然后结束。它会先在离子静止系内解出：

\begin{itemize}
\item 更新后的 electron
\item 更新后的 ion
\item photon 三动量
\end{itemize}

再一起变回 lab frame。最后写到 photon species 的不是“真实质量意义下的速度”，而是

\[
u = \frac{p}{m_e}
\]

型的统一接口变量：

% <-- fenced code (cpp)
\begin{codeblock}
upx = p3x_rel / PhysConst::m_e;
upy = p3y_rel / PhysConst::m_e;
upz = p3z_rel / PhysConst::m_e;
\end{codeblock}

这样做不是修改 photon 物理，而是为了让 collision-created photons 继续复用 WarpX 现有的粒子 SoA 与 diagnostics 链。

对 fusion 来说，上一节已经区分了 D-D / D-T 两产物和 p-B11 三 alpha 两条物理路径；这一层再补的是“为什么实现里看到的宏粒子数更多”。两产物 fusion 的 \cpath{TwoProductFusionInitializeMomentum(...)} 其实只算一次 products 动量，然后把：

\begin{itemize}
\item 第一 product 的一组 \code{(ux,uy,uz)} 写两次
\item 第二 product 的一组 \code{(ux,uy,uz)} 也写两次
\end{itemize}

所以实现里出现 4 个宏粒子，并不表示一次 fusion 物理上有 4 个独立 products，而是因为 WarpX 在两个 parent 位置各复制一套 child。

p-B11 的 \cpath{ProtonBoronFusionInitializeMomentum(...)} 进一步复杂一些。它先按

% <-- fenced code (text)
\begin{codeblock}
p + B11 -> alpha1 + Be8*
\end{codeblock}

做一次两产物动量求解，再在 \cpath{Be8*} 静止系里把剩余 \code{3.12 MeV} 衰变能随机分到：

% <-- fenced code (text)
\begin{codeblock}
Be8* -> alpha2 + alpha3
\end{codeblock}

最后再变回 lab frame。于是实现里会看到 6 个 alpha 宏粒子：

\begin{itemize}
\item \cpath{alpha1} 两份
\item \cpath{alpha2} 两份
\item \cpath{alpha3} 两份
\end{itemize}

本质上仍然是“在两个 parent 位置各复制一套 products”的事件记账方式，而不是 6 个不同物理 alpha 通道。

所以到这一层可以把 collision 剩余 kernel 细节压成三句：

\begin{enumerate}
\item Perez 路径最关键的是局域尺度闭合和加权 \cpath{n12}
\item Bremsstrahlung 最关键的是 photon 动量写回并与 parent 更新同时完成
\item Fusion 最关键的是“先算真实 kinematics，再做双 parent 位置复制”的宏粒子实现语义
\end{enumerate}

\subsection{更深一层的 event kernel：Perez 角分布、Bremsstrahlung 能谱反演与 fusion 概率控制}

如果再往 \cpath{BinaryCollision} 的 event kernel 里读，collision 还有一层之前没写清的共同问题：上一节讲的是“哪些量会被写回”，这一层讲的是“这些量到底怎样被抽出来”。

对 \cpath{pairwisecoulomb} 而言，\cpath{ElasticCollisionPerez(...)} 上游已经准备好了：

\begin{itemize}
\item \cpath{bmax}
\item \cpath{sigma\_max}
\item 加权后的 \cpath{n12}
\end{itemize}

但 \cpath{UpdateMomentumPerezElastic(...)} 并不是把这些量直接当成一次 Bernoulli 碰撞概率。它先在 COM frame 内构造归一化散射长度
\[
s_{12},
\]
然后按 \cpath{s12} 所处区间切到四段式角分布：

\begin{itemize}
\item \code{s12 <= 0.1}：\code{cosXs = 1 + s12 * log(r)}
\item \code{0.1 < s12 <= 3}：用五次多项式近似的 \cpath{Ainv}
\item \code{3 < s12 <= 6}：用 \code{A = 3 exp(-s12)} 的另一段反演
\item \code{s12 > 6}：直接退化成 \code{cosXs = 2r - 1} 的各向同性散射
\end{itemize}

随后再抽一个独立方位角，把 post-collision 动量从 COM frame 变回 lab frame。更关键的是，最后并不是无条件同时更新两只 reactant，而是分别按：

% <-- fenced code (cpp)
\begin{codeblock}
if (w2 > r * max(w1, w2)) update particle 1
if (w1 > r * max(w1, w2)) update particle 2
\end{codeblock}

做 weighted rejection。因此 WarpX 当前的 weighted macroparticle Coulomb collision，真正的语义是：

\begin{itemize}
\item \cpath{n12} 和 \cpath{s12} 决定散射强度
\item reactant 是否真正写回，再由宏粒子权重比决定
\end{itemize}

对 \cpath{bremsstrahlung} 而言，前两节已经区分了 \cpath{BremsstrahlungFunc} 和 \cpath{PhotonCreationFunc}。更深一层看 \cpath{BremsstrahlungEvent(...)}，会发现它的前半段先把电子变到离子静止系，再结合电子密度构造 plasma-frequency cutoff：
\[
E_{\mathrm{cut}} = \hbar \omega_{pe}.
\]
如果电子在离子静止系里的动能 \cpath{KE\_eV} 低于这个阈值，就直接返回 0，不让软光子进入采样。若允许发生事件，代码并不是从表里拿一个最近点，而是先按电子能量方向插值，再在 \cpath{k/T1} 网格上逐段积累 trapezoidal cross section，最后在命中的区间内反演出连续 photon energy。

这一层还有一个必须记录的当前实现边界：event probability 用的核心量是
\[
\mathrm{arg}
=
f_{\mathrm{multi}}\, v_1\, \sigma_{\mathrm{total}}\, n_2\, dt\,
\frac{\gamma_1^{\mathrm{rel}}}{\gamma_1 \gamma_2}.
\]
当 \code{arg > 1} 时，源码会先把 \cpath{arg} 饱和到 1，再去改 \cpath{fmulti}。按当前实现，这一分支直接生效的是 probability saturation，而不是像 fusion 那样形成真正有效的 multiplier backoff。这一点应当被视为当前源码边界，而不是已有功能。

\cpath{nuclearfusion} 的 \cpath{SingleNuclearFusionEvent.H} 则正好相反：这里的 probability control 是真的做完了。它先算
\[
P_{\mathrm{est}}
=
\mathrm{multiplier\_ratio}\,
f_{\mathrm{mult}}\,
\mathrm{lab\_to\_COM\_factor}\,
w_{\max}\,
\sigma_f(E_{\mathrm{coll}})\,
v_{\mathrm{coll}}\,
\frac{dt}{dV},
\]
如果超过 \cpath{probability\_threshold}，就按 \cpath{probability\_target\_value} 回退到一个新的 \code{fusion_multiplier_eff >= 1}，再把 product weight 缩成
\[
w_{\mathrm{product}} = \frac{w_{\min}}{f_{\mathrm{mult,eff}}}.
\]
最终真正使用的事件率也不是简单线性截断，而是
\[
P = 1 - e^{-P_{\mathrm{est}}},
\]
在代码里写成 \cpath{-std::expm1(-probability\_estimate)}，专门保证小概率时的数值稳定性。

这一层和 regression 的关系也更清楚了：

\begin{itemize}
\item 2D / 3D Coulomb tests 看的主要是指数 relaxation fit
\item \codeesc{inputs\_test\_3d\_collision\_iso\{,\_subcycle\}} 看的是各向异性温度的 isotropization 解析解，以及 \cpath{ndt\_subcycle} 是否保持同一 collision timestep
\item \cpath{inputs\_test\_rz\_collision} 看的是 cylindrical-cell 配对与临时动量旋转假设下，局域共线粒子几乎不发生有效散射
\item \cpath{background\_mcc} 的 \cpath{capacitive\_discharge} 条目则不该再混成普通 collision 单元测试：1D PICMI 入口实际拆成 \cpath{analysis\_1d.py} 与 \cpath{analysis\_dsmc.py} 两条 case-1 profile 对照，前者验证 \code{background_mcc + external Poisson solver callback}，后者验证把 DSMC 分支接入同一骨架后的离子密度 profile；2D native / PICMI 入口当前仍主要是应用级 checksum 基线，而 \cpath{test\_2d\_background\_mcc\_dp\_psp} 在当前 \cpath{CMakeLists.txt} 中仍是注释掉的遗留变体
\end{itemize}

\subsection{性能与数值后处理层：resampling、thermalizer 与 sorting}

在 collision/QED 这些多物理分叉之外，\cpath{Particles/} 里还有一层更偏数值控制与性能优化的对象图：

\begin{itemize}
\item \cpath{Resampling}
\item \cpath{ParticleThermalizer}
\item \cpath{Sorting}
\end{itemize}

它们不直接改变场求解器，但会显著改变宏粒子数、粒子分布以及后续 kernel 的访存局部性。

\cpath{Resampling} 是 species 级开关。\cpath{PhysicalParticleContainer.cpp} 先读 \cpath{do\_resampling}，再按 species 构造一个 \cpath{Resampling} 对象。它内部又拆成：

\begin{itemize}
\item \cpath{ResamplingTrigger}
\item \cpath{ResamplingAlgorithm}
\end{itemize}

触发条件不是只有固定步数，而是
\[
\texttt{intervals.contains(step)}
\;\lor\;
\left(\frac{N_{\mathrm{global}}}{N_{\mathrm{cells,global}}} >
\texttt{resampling\_trigger\_max\_avg\_ppc}\right).
\]

调用位置在主循环里对应的是 \cpath{istep[0]+1}，所以匹配的是“下一步编号”。一旦触发，WarpX 会先 \cpath{Redistribute()}，再在各 level/tile 上调用具体算法，最后统一 \cpath{deleteInvalidParticles()}。因此 resampling 和 collision、absorbing boundary、EB scraping 一样，仍沿用“kernel 内先标 invalid，之后统一删粒子”的合同。

当前正式接入的算法有两条。\cpath{LevelingThinning} 是按 cell 独立工作的：对每个 cell 先算平均权重
\[
\bar w_{\mathrm{cell}},
\]
再定义
\[
w_{\mathrm{level}} = \bar w_{\mathrm{cell}} \times \texttt{target\_ratio}.
\]
凡是 \code{w_i <= w_level} 的粒子，就以
\[
1 - \frac{w_i}{w_{\mathrm{level}}}
\]
的概率删除，否则把权重直接抬到 \cpath{w\_level}。这条算法只会抬低权重粒子，不会动大权重尾部。对应的 \cpath{test\_2d\_leveling\_thinning} analysis 也比较强：一类 species 检查最终粒子数和统一权重，一类 species 检查 Gaussian 权重分布下的 level-weight 和 untouched heavy tail。

\cpath{VelocityCoincidenceThinning} 则不是简单抽稀，而是 merge。它先在每个 cell 内按速度空间分 bin，可以是 spherical grid，也可以是 cartesian grid；然后把同一个 velocity bin 内的 cluster 压成两个保留粒子。实现上会累计 cluster 的总权重、加权平均位置、加权平均动量和总动能，再构造一对关于平均动量对称的新动量，使 cluster 内的线动量和动能都守恒，其余粒子全部标 invalid。这里 \cpath{resampling\_algorithm\_target\_weight} 的实际语义也要记清：内部会乘 2，因为每个 cluster 最后固定压成两粒子。当前这条路径只有 checksum，没有像 \cpath{LevelingThinning} 那样的独立 analysis。

\cpath{ParticleThermalizer} 则是单独的全局参数块 \cpath{particle\_thermalizer.*}，不放在 species 段里。它当前只支持 Cartesian 1D/2D/3D，不支持 \cpath{RZ/RCYLINDER/RSPHERE}。在主循环中的插入位置非常具体：

\begin{enumerate}
\item moving window
\item particle boundary / EB scraping / invalid 删除 / sorting
\item \cpath{m\_particle\_thermalizer.applyThermalizer(*mypc)}
\item collisions
\end{enumerate}

因此 thermalizer 的效果会先落到“已经通过边界筛选保留下来的粒子”上，再进入碰撞模块。它也不是硬墙式 momentum reset，而是在一个有法向、起止位置和渐进概率的 thermal region 内，对超过 \cpath{momentum\_threshold} 的动量分量抽样重置为方差由 \cpath{theta} 决定的 Gaussian。\cpath{particle\_absorbing\_boundary} 会显式打开 \cpath{particle\_thermalizer.*}，并用 \cpath{PhaseSpaceElectrons} 直方图断言吸收边界附近的反向高速电子权重显著降低；但这仍是 absorbing boundary、field-function laser、reduced diagnostic 和 thermalizer 的耦合 regression，不是 dedicated thermalizer-only 单测。

\cpath{Sorting} 则必须和前文 AMR coarse-fine 的 \cpath{PartitionParticlesInBuffers()} 区分开来。后者是物理分流，前者是全局性能阶段。\cpath{WarpXEvolve.cpp} 在 boundary 处理之后检查 \cpath{sort\_intervals.contains(step+1)}，触发后调用：

% <-- fenced code (cpp)
\begin{codeblock}
mypc->SortParticlesByBin(
    sort_bin_size, m_sort_particles_for_deposition, m_sort_idx_type);
\end{codeblock}

文档和 \cpath{WarpX.cpp} 一起看，当前默认值是平台相关的：

\begin{itemize}
\item GPU 默认 \code{sort_intervals = 4}
\item CPU 默认 \code{sort_intervals = -1}
\end{itemize}

目的是提升 memory locality，不是物理修正。并且还有两种模式：

\begin{itemize}
\item \code{sort_particles_for_deposition = true}
\begin{itemize}
\item 走 deposition-specialized sort
\end{itemize}
\item \code{sort_particles_for_deposition = false}
\begin{itemize}
\item 走 generic bin sort，粒度由 \cpath{sort\_bin\_size} 控制
\end{itemize}
\end{itemize}

所以到这里，\cpath{Particles/} 的“非主物理 kernel 层”也形成了清晰分层：

\begin{itemize}
\item resampling：调粒子数与权重/局域 phase-space 代表性
\item thermalizer：调局部高动量粒子的热化行为
\item sorting：调后续 deposition/gather 的访存局部性
\end{itemize}

\subsection{Vranic 2015：两粒子 merge 的论文-源码边界}

Vranic 等人的论文把粒子合并明确放在六维 phase space 中：先用空间 merge cell 和动量 cell 找到相近 cluster，再把一个 cluster 压成两个新宏粒子。单个新粒子一般不能同时满足总权重、总动量、总能量和质量壳关系；两粒子构造则可以令

\[
w_t=w_a+w_b,\qquad
\boldsymbol{p}_t=w_a\boldsymbol{p}_a+w_b\boldsymbol{p}_b,\qquad
\epsilon_t=w_a\epsilon_a+w_b\epsilon_b.
\]

在等权、等能量的简化下，两个新动量关于总动量方向对称，论文用

\[
\cos\theta=\frac{p_t}{w_t p_a}
\]

确定夹角，并用动量 cell 对角线选择平面，避免在原来没有展宽的方向凭空制造分布。论文还强调位置不能机械地都放在 cluster 质心，而应从原粒子位置中抽取，以避免 merge cell 中心出现人为密度尖峰。论文的图和数值例子说明了这种构造怎样在具体分布中工作；阅读时应把它们视为合并算法的机制证据，而不是 WarpX 默认设置的性能承诺。

这与 WarpX 的 \cpath{VelocityCoincidenceThinning} 存在清晰的结构映射：源码同样在 cell 内按速度空间分 bin，把 cluster 压成两个粒子，并累计权重、加权动量和动能；\cpath{resampling\_algorithm\_target\_weight} 还因“两粒子输出”而在内部乘 2。论文的 Poisson merging-rate 公式、QED cascade speed-up 和分布复现结果则属于论文自己的 OSIRIS/QED 案例，不能转写成 WarpX 已完成的 runtime physics proof。

WarpX 的 \cpath{resampling} regression 主要检查粒子数、权重或 checksum；它没有把论文中的 two-stream、magnetic-shower 和 QED-cascade 逐一搬到相同的输入和诊断上。因此这里能成立的结论是“论文方法与 WarpX 的结构可以对应”，而不是“WarpX 已复现论文算例”或“加速比已被证明”。要评估一个重采样设置，至少还应比较重采样前后的局部总权重、动量、能量、空间分布和权重尾，而不能只看最终粒子数。

\subsubsection{Muraviev 2021：agnostic down-sampling 的论文-源码边界}

Vranic 2015 解释了“一个 cluster 压成两个粒子时如何同时保持局部动量和能量”；Muraviev 等人进一步把重采样问题拆成 merging、thinning 和 complete resampling，并提出 agnostic down-sampling 原则：至少一个粒子的权重变为零，同时每个原粒子的期望新权重仍等于旧权重。由此得到的不是单次 realization 的严格局部不变，而是任意由位置、动量或其他粒子状态定义的分布在 ensemble average 下保持不变。

论文比较了 \cpath{simple}、\cpath{leveling}、\cpath{globalLev}、\cpath{numberT}、\cpath{energyT}、\cpath{conserv}、\cpath{mergeAv} 和 \cpath{merge}。其中 \cpath{numberT} 严格保持 cell 总权重，\cpath{energyT} 严格保持 cell 总能量，\cpath{conserv} 可以把能量、三分量动量、总权重和空间中心矩组成线性不变量；反过来，\cpath{simple} 虽然最容易实现，却会产生很宽的局部权重尾，在 QED cascade 中可能制造无法被时间步解析的局部等离子体频率和非物理场增长。

这篇论文与 WarpX 有三条可用的概念连接：一是 \cpath{LevelingThinning} 对低权重粒子的 leveling 思路；二是 \cpath{VelocityCoincidenceThinning} 在 velocity bin 内把 cluster 压成两个粒子的 merge 结构；三是“只看 checksum 不足以证明重采样物理质量”的验证要求。论文的 PICADOR/hi-chi 运行使用了自己的 QED cascade、Weibel 和 k-means 实验，不能把其 growth rate、运行时间、权重尾或图 1--12 数值直接写成 WarpX 结果。

对 WarpX 而言，这篇论文最重要的启发是验证标准：算法名称、粒子数下降或 checksum 一致都不足以说明分布质量。若要判断一次 thinning 或 merge 是否可接受，应在同一空间单元和动量区间比较局部总权重、能量和三分量动量，另外报告 density variance 与权重尾；这些量分别对应守恒、统计代表性和数值稳定性。论文中 PICADOR/hi-chi 的 QED cascade、Weibel 等结果属于其自身的物理设置，不能直接充当 WarpX 的运行证据。

\subsection{四类单粒子问题：该测什么，不能推出什么}

\cpath{particle\_pusher}、\cpath{single\_particle}、\cpath{larmor} 和 \cpath{photon\_pusher} 都只有很少的粒子，却不是同一种“单粒子测试”：

\begin{itemize}
\item \cpath{particle\_pusher}
\item \cpath{single\_particle}
\item \cpath{larmor}
\item \cpath{photon\_pusher}
\end{itemize}

读者首先要问的不是粒子数，而是每个输入固定了哪一种物理情形、比较了哪个观察量，以及该比较没有覆盖什么。

\cpath{particle\_pusher} 是最直接的一条。输入里固定：

% <-- fenced code (text)
\begin{codeblock}
algo.particle_pusher = "higuera"
particles.B_ext_particle_init_style = "constant"
particles.B_external_particle =  0.0  0.0  1.0
particles.E_ext_particle_init_style = "constant"
particles.E_external_particle =  -2.994174829214179e+08  0.0  0.0
\end{codeblock}

并让单个 positron 在满足
\[
E_x = -v_y B_z
\]
的 force-free 场中跑 \cpath{10000} 步。\cpath{analysis.py} 只检查最终
\[
x \approx 0.
\]
因此它并不是一般轨道画图，而是直接给：

\begin{itemize}
\item \cpath{PushSelector.H} 的 \cpath{ParticlePusherAlgo::HigueraCary}
\item \cpath{UpdateMomentumHigueraCary(...)}
\item \cpath{UpdatePosition(...)}
\end{itemize}

这整条 relativistic Higuera-Cary push 主链做强断言。

该输入得到的

\[
\max |x|=1.1430664324\times10^{-4},
\]

低于 \codeesc{10\textasciicircum{}\{-3\}} 容差。它支持“此 force-free 条件下的相对论推进和位置更新彼此相容”，但只覆盖单粒子、恒定外场和横向偏移这一观察量，不能替代三种推进器的完整轨道、能量或相空间基准。

在同一输入中只替换 \cpath{algo.particle\_pusher}，可得到一个有意窄化的算法对照：

% <-- table block (source: 04-particle-pushers L1828)
\begin{tabularx}{\textwidth}{LrrL}
\toprule
pusher & 末态 $\max|x|$ & \cpath{1e-3} gate & 解释 \\
\midrule
\endfirsthead
\toprule
pusher & 末态 $\max|x|$ & \cpath{1e-3} gate & 解释 \\
\midrule
\endhead
Boris & \cpath{2.3213958529e3} & FAIL & force-free cancellation 在这个高相对论设置下明显失真 \\
Vay & \cpath{1.0795497978e-4} & PASS & 保留较好的 relativistic frame/cancellation 行为 \\
Higuera--Cary & \cpath{1.1430664324e-4} & PASS & 在该观察量上与 Vay 同量级 \\
\bottomrule
\end{tabularx}

三组不是三份独立的官方输入，而是对同一个 force-free 问题的 pusher-only 对照。因此它只支持相对论 cancellation 的差异提示，不替代 boosted-frame、Poincare section 或长期能量/相空间 benchmark。

\cpath{single\_particle} 则必须拆成两类。

第一类 \cpath{inputs\_test\_2d\_bilinear\_filter} 虽然也只有一个电子，但 analysis 完全不看轨道，而是手工构造未滤波 \cpath{Jx}、再用二维 bilinear kernel 卷积，最后和 plotfile 里的 \cpath{jx} 比较。它真正验证的是：

\begin{itemize}
\item 单粒子沉积后的 current stencil
\item \cpath{warpx.use\_filter}
\item \cpath{warpx.filter\_npass\_each\_dir}
\item 对称边界 bilinear filtering 的实现
\end{itemize}

所以它更接近 deposition/filter regression。

第二类 \cpath{inputs\_test\_1d\_synchronize\_velocity} 则在 \code{algo.maxwell_solver = none} 下给一个电子施加常量 \cpath{E\_z}，同时打开：

% <-- fenced code (text)
\begin{codeblock}
warpx.synchronize_velocity_for_diagnostics = 1
\end{codeblock}

analysis 先在 Python 里手动做 half-backward、5 步 leapfrog、再加 half-forward 得到同步后的速度，然后和 diagnostics 输出比较。这条 regression 直接对应：

\begin{itemize}
\item 参数 \cpath{warpx.synchronize\_velocity\_for\_diagnostics}
\item \cpath{WarpXEvolve.cpp} 里 diagnostics 前的 \cpath{SynchronizeVelocityWithPosition()}
\end{itemize}

也就是说，它验证的不是 Boris 物理轨道误差，而是“输出给 diagnostics 的速度是否和位置处在同一时间层”。

第 5 个诊断步的模拟/理论值分别为：

\begin{itemize}
\item $z=2.2985203786002786/2.2985203756075920$；
\item $u_z=879410.0053860814/879410.0053860815$；
\item 相对速度误差为 \code{1.3237889e-16 < 1e-15}。
\end{itemize}

这支持的是 diagnostics time-level synchronization，不应被误写成单粒子 pusher 的独立轨道精度 benchmark。

\cpath{photon\_pusher} 又是另一类。输入里建了 16 个 photon species，覆盖：

\begin{itemize}
\item 六个坐标轴正负方向
\item 对角方向
\item 动量大小 \cpath{1} 和 \cpath{10}
\end{itemize}

analysis 用理论式
\[
\mathbf{x}(t)=\mathbf{x}_0 + ct\,\hat{\mathbf{u}},
\qquad
\mathbf{p}(t)=\mathbf{p}_0
\]
逐 species 比较最终位置和动量。因此它真正打到的是：

\begin{itemize}
\item \cpath{PhotonParticleContainer::PushPX()}
\item \cpath{UpdatePosition(...)} 里的 massless branch
\item photon 不做 charge/current deposition 的合同
\end{itemize}

16 个 photon species 的末态最大相对误差为：

\begin{itemize}
\item 位置直线传播：\code{6.0986372e-16 < 1e-14}；
\item 动量保持：\code{1.7217530e-16 < 2.2204460e-16}。
\end{itemize}

这条证据支持的是无质量粒子的 \cpath{c} 速率传播、方向保持和不参与带电粒子 current deposition 的路径，不应与 Boris/Vay/Higuera--Cary 的带电动量旋转混写。

最后 \cpath{larmor} 反而要最保守。它输入里组合了：

\begin{itemize}
\item \cpath{electron} 和 \cpath{positron}
\item 常量外部粒子磁场 \cpath{B\_y}
\item \code{amr.max_level = 1}
\item PML
\item \code{warpx.do_dive_cleaning = 1}
\item full/raw diagnostics
\end{itemize}

但 \cpath{CMakeLists.txt} 里没有独立 analysis，只有 checksum。因此更准确的说法是：

\begin{itemize}
\item 这是 charged-particle gyro-motion 与 external-particle-field、MR、PML、div-cleaning 组合稳定性的 checksum 基线；
\item 它不是已有独立解析半径/回旋频率对照的强单粒子 analysis。
\end{itemize}

若把这个组合输入直接同连续 uniform-\cpath{B\_y} 轨道比较，电子和正电子的轨迹相对位移误差都是 \cpath{1.28285096e-2}，动量相对误差都是 \cpath{9.69641193e-2}。这不否定 checksum 的作用；它只说明 MR/PML/divergence-cleaning 的组合已经超出连续单粒子解析解的假设。因此本章把它保留为 checksum-only 基线，不能用它升级为强物理 gate。

这四类输入共同给出一个实用的阅读法：

\begin{itemize}
\item Higuera--Cary force-free relativistic push：检查相对论推进链。
\item diagnostics 半步速度同步：检查输出时间层。
\item bilinear current filter：检查沉积后的网格量。
\item photon 直线传播与动量守恒：检查无质量传播。
\item Larmor 半径/频率独立解析对照：这组输入没有提供，不能夸大。
\item Larmor 连续轨道比较：只用于说明为何不能把该组合输入提升为强 gate。
\end{itemize}

把它们都叫作“单粒子测试”，会掩盖观察量、源码路径和可支持结论之间最关键的差别。

\subsection{推进器修改后的验证阶梯：先选对 consumer，再解释结果}

当读者修改 pusher、外部粒子场、位置更新或 diagnostics 选项时，最常见的错误不是没有输出，而是把一个通过的单粒子测试扩张成了不属于它的结论。下面四层不是按“测试难度”排序，而是按被 consumer 实际读取的状态排序。

\textbf{第一层：带质量粒子的 momentum--position 链。}若改动的是 \cpath{ParticlePusherAlgo::HigueraCary} 分派、\cpath{UpdateMomentumHigueraCary()} 或 \cpath{UpdatePosition()} 的显式带质量路径，先回到 \cpath{test\_3d\_particle\_pusher}。它固定一个 \cpath{SingleParticle} positron、\code{algo.particle_pusher = "higuera"}、常量粒子侧 \cpath{E\_x/B\_z}、\code{max_step = 10000}，并由 CTest 把 \cpath{diags/diag1010000} 交给 \cpath{analysis.py}。consumer 只读取末态 \cpath{particle\_position\_x}，断言 \code{abs(x) < 1e-3}。这个量对应 force-free cancellation：在该输入下 (E\_x=-v\_yB\_z)，横向位置不应出现大的漂移。

通过这层只支持一个窄结论：给定这一个外场、初始动量、10000 步和 Higuera--Cary 分支，动量更新与位置更新保持该 force-free observable。它不能证明 Boris、Vay、任意电磁场、particle deposition、AMR 或自洽场都正确。若只把同一输入的 \cpath{algo.particle\_pusher} 改成另一种算法，原来的 (x) consumer 仍可作为同一 force-free 问题的比较量，但结果必须标为对照实验，不能把它改写成该算法的官方通用认证。

\textbf{第二层：输出时间层，而不是轨道算法。}若改动的是 \cpath{warpx.synchronize\_velocity\_for\_diagnostics}、diagnostic 写出时机或速度同步代码，应使用 \cpath{test\_1d\_synchronize\_velocity}。它以常量 (E\_z) 推进一个电子，在第 5 步写 Full diagnostics；\cpath{analysis\_synchronize\_velocity.py} 从 half-backward、五次 leapfrog 与 half-forward 显式重建同步后的 (u\_z)，并比较 \cpath{diags/diag1000005} 中的位置与动量，速度相对误差阈值为 \cpath{1e-15}。即使它通过，也只能说明 diagnostics 读到的位置与速度时间层相容；它不能证明 Higuera--Cary、Boris 或 Vay 的相对论轨道精度。

\clearpage

\textbf{第三层：无质量粒子是另一条容器链。}若改动 \cpath{PhotonParticleContainer::PushPX()}、massless 的 \cpath{UpdatePosition()} 分支，使用 2-rank \cpath{test\_3d\_photon\_pusher}。输入为 16 个 photon species，覆盖轴向/对角方向和两档动量；analysis 在 \cpath{diags/diag1000050} 中逐 species 读取末态。

它分别对照直线位置与初始动量：

\[
\mathbf{x}(t)=\mathbf{x}_0+ct\,\hat{\mathbf u},
\qquad
\mathbf p(t)=\mathbf p_0.
\]

它的 position 与 momentum consumer 分别施加 \cpath{1e-14} 和 machine-epsilon 阈值。源码中的 photon 容器在 \cpath{PositionPushType::Full} 时调用同一个 \cpath{UpdatePosition()}，但又因 photon 不带电而跳过 current deposition；因此它不能替代带质量 pusher 的 Lorentz force 验证，也不能替代第 5 章的 charge/current 合同。

\textbf{第四层：checksum 仍有价值，但不是解析 gate。}\cpath{test\_2d\_larmor} 组合了外部粒子磁场、两层网格、PML 与 divergence cleaning，却在 CMake 中把独立 analysis 标为 \cpath{OFF}，只保留末态 checksum。它适合发现这组组合输入的输出回归，却没有提供独立半径或回旋频率 consumer。不要因为输入只有电子和正电子，就把 checksum 通过写成“Larmor 轨道已解析验证”。

实际排错可以遵循一个简短顺序：先问被改的是带质量 momentum/position、diagnostic time level、massless position，还是 deposition/field；只对前三类分别选上面对应的 consumer。若 force-free (x) 失败，检查 pusher 选择、常量 external particle fields、初始 (u\_y)、步数和末态 diagnostic；若只有 synchronized (u\_z) 失败，优先检查 diagnostics 前的速度同步；若 photon 的位置或动量失败，转查 photon container 和 massless position branch。若改动触及 charge/current、field solver 或 AMR，单粒子通过不能完成验证，应转入第 5、6、7 章对应的 source、场和边界 consumer。

\subsection{粒子诊断与外场：两条不经过主网格场的路径}

在“单粒子/推进器”之外，\cpath{particle\_fields\_diags} 与 \cpath{plasma\_lens} 分别回答两个不同的问题：怎样把粒子属性归约为网格诊断量，以及怎样在不读取主网格场的条件下给粒子施加外场。

\begin{itemize}
\item \cpath{particle\_fields\_diags}
\item \cpath{plasma\_lens}
\end{itemize}

它们都直接消费粒子位置、动量或外部粒子场，但验证对象并不是普通 pusher 轨道。

\cpath{particle\_fields\_diags} 的关键输入不是单粒子初值，而是 diagnostics 配置：

% <-- fenced code (text)
\begin{codeblock}
diag1.particle_fields_to_plot = z uz uz_filt zuz jz
diag1.particle_fields_species = electrons protons photons
diag1.particle_fields.z(x,y,z,ux,uy,uz) = z
diag1.particle_fields.uz(x,y,z,ux,uy,uz) = uz
diag1.particle_fields.uz_filt(x,y,z,ux,uy,uz) = uz
diag1.particle_fields.uz_filt.filter(x,y,z,ux,uy,uz) = (uz < 0)
diag1.particle_fields.zuz(x,y,z,ux,uy,uz) = z * uz
diag1.particle_fields.jz(x,y,z,ux,uy,uz) = uz*q_e
diag1.particle_fields.jz.do_average = 0
\end{codeblock}

analysis 会同时做三件事：

\begin{enumerate}
\item 用 \cpath{yt} 从粒子数据直接手工重建 cell-centered quantity
\item 读取 Full plotfile 中的 \cpath{boxlib} particle-field meshes
\item 再读取 openPMD 中同名 meshes
\end{enumerate}

然后逐一比较：

\begin{itemize}
\item \cpath{zavg}
\item \cpath{uzavg}
\item \cpath{zuzavg}
\item \cpath{uzavg\_filt}
\item \cpath{jz}
\end{itemize}

因此它真正验证的是：

\begin{itemize}
\item \cpath{Diagnostics.cpp} 对 \cpath{particle\_fields\_to\_plot}、\cpath{.filter(...)}、\cpath{.do\_average} 的参数解析
\item \cpath{ParticleReductionFunctor} 如何把粒子归约成 cell-centered diagnostics field
\item plotfile 与 openPMD writer 在这一层是否一致
\end{itemize}

这说明 WarpX 不仅能把粒子作为散点写出去，还能按用户给出的 parser 表达式把 species 内的粒子数据重新投影成网格诊断量。

\cpath{plasma\_lens} 则落在粒子侧外场主链上。analysis 不是读主网格场，而是直接比较两颗测试电子穿过 lens 序列后的：

\begin{itemize}
\item 最终横向位置
\item 最终横向动量
\end{itemize}

与解析透镜串联模型是否一致。

这里要区分两条源码入口。

第一条是：

% <-- fenced code (text)
\begin{codeblock}
particles.E_ext_particle_init_style = repeated_plasma_lens
particles.B_ext_particle_init_style = repeated_plasma_lens
\end{codeblock}

它对应：

\begin{itemize}
\item \cpath{MultiParticleContainer.cpp} 读取 lens 周期、起点、长度、强度
\item \cpath{GetExternalEBField} 在粒子 gather 之后、push 之前按粒子位置实时叠加 lens focusing field
\end{itemize}

第二条是 hard-edged 版本：

% <-- fenced code (text)
\begin{codeblock}
lattice.elements = ...
plasmalens*.type = plasmalens
plasmalens*.dEdx = ...
\end{codeblock}

它对应 accelerator lattice 分支里的 \cpath{HardEdgedPlasmaLens}。

两者共用同一个 analysis，因此真正被验证的是：

\begin{itemize}
\item repeated-plasma-lens 粒子侧外场
\item accelerator-lattice hard-edged lens
\item boosted-frame plasma-lens 反变换后一致性
\item short-lens residence correction
\end{itemize}

这一组的教学意义不是“又一条轨道”，而是两条与普通自洽场推进不同的路径：

\begin{itemize}
\item 粒子 diagnostics 可以把粒子属性重新压成 cell-centered field
\item 粒子侧外场可以完全绕过主网格场寄存器，直接通过 \cpath{GetExternalEBField} 进入 \cpath{PushPX()}
\end{itemize}

accelerator lattice 还有一个直接的解析基准：\cpath{Examples/Tests/accelerator\_lattice/hard\_edged\_quadrupoles*}。单电子穿过 \code{drift + quad + drift + quad} 串联，分析从输入重建 \cpath{lattice.elements}、\cpath{drift.ds}、\cpath{quad.ds}、\cpath{quad.dEdx}，再按 hard-edged quadrupole 的解析透镜公式逐段积分，要求最终 \cpath{x} 误差低于 \codeesc{1\%}、\cpath{u\_x} 误差低于 \codeesc{0.2\%}。boosted-frame 与 moving-window 变体继续使用同一解析对照，因此这里检查的不是“lattice 参数能读入”，而是 \cpath{HardEdgedQuadrupole}、\cpath{LatticeElementFinder} 与 \cpath{PushPX()} 能否共同给出正确的粒子偏转。

这里还有一个必须分清的源码边界：\cpath{drift} 只把长度 \cpath{ds} 转成 beamline 的 \cpath{zs/ze} 几何区间，不直接返回外场。真正给粒子累加 \cpath{E/B} 的是 \cpath{HardEdgedQuadrupole} 与 \cpath{HardEdgedPlasmaLens}；\cpath{LatticeElementFinder} 按 tile 建立最近元件索引，把 boosted-frame 粒子坐标和步末 \cpath{z+v\_zdt} 变回实验室系，调用各元件的 \cpath{get\_field(...)}，再把累计场变回 boosted frame 后交给 \cpath{PushPX()}。所以 \code{drift + quad + drift + quad} 中的 drift 只影响解析 beamline 几何和 residence 区间判定，不参与外场累加。

读这一节时可用两个问题自检：写出的 particle-field mesh 是否等于从粒子数据重新归约的量？粒子所见的透镜场是否来自外场元件，而不是主网格场？前者要求同时比较数据归约和 writer，后者要求同时比较元件几何、参考系变换和位置/动量的解析结果。

\subsection{边界缓冲区与 Python 粒子操作：能控制什么，尚未覆盖什么}

\cpath{particle\_boundary\_scrape}、\cpath{particle\_data\_python} 和 single-precision particle fields 都不以轨道为主要观察量。它们分别回答：粒子被边界删除后在哪里可见，Python 能否读写粒子状态并触发沉积，以及单精度归约是否已有可执行的误差检查。

\begin{itemize}
\item \cpath{particle\_boundary\_scrape}
\item \cpath{particle\_data\_python}
\item \cpath{particle\_fields\_diags} single-precision FIXME
\end{itemize}

阅读这些输入时，应把验证对象分成三层：

\begin{itemize}
\item scraped-particle buffer
\item Python runtime attribute / injection / deposition wrapper
\item 单精度粒子 diagnostics 误差边界
\end{itemize}

\subsubsection{被删除的粒子：主容器与边界缓冲区必须同时检查}

\cpath{particle\_boundary\_scrape} 的 native 输入配置一个立方体 EB 和一束电子，\cpath{analysis\_scrape.py} 做两个直接检查：

\begin{itemize}
\item 第 40 步还应有 612 个电子
\item 第 60 步主 species 中电子数应变成 0
\end{itemize}

这说明 \cpath{ScrapeParticlesAtEB()} 确实把撞到 embedded boundary 的粒子删掉了。

PICMI 变体在 \cpath{sim.step(...)} 后构造 \cpath{ParticleBoundaryBufferWrapper()}，并检查：

\begin{itemize}
\item EB buffer 中累计粒子数是 612
\item \cpath{stepScraped} 全都大于 40
\item 所有 rank 汇总后的 buffer 粒子总数仍是 612
\item \cpath{clear\_buffer()} 之后 buffer size 回到 0
\end{itemize}

因此读者应同时检查两个对象：

\begin{enumerate}
\item EB scraping 确实把粒子从主容器里删除
\item 删除掉的粒子确实进入了 Python 可访问、可清空的 boundary buffer
\end{enumerate}

\subsubsection{Python 操作粒子：属性、注入与沉积是三条独立操作}

\cpath{particle\_data\_python} 没有独立 \cpath{analysis.py}，断言直接写在 PICMI 输入脚本里，主要步骤是：

\begin{itemize}
\item \cpath{sim.initialize\_warpx()}
\item \code{sim.particles.get("electrons")}
\item \code{add_real_comp("newPid")}
\item 在 \cpath{beforestep} callback 里持续 \cpath{add\_particles(...)}
\item 再直接断言 \cpath{get\_real\_comp\_index(...)}、tile 里的 \cpath{newPid} 值，以及 Python wrapper 暴露的 \cpath{deposit\_current(...)} 确实能把电流沉到 \cpath{current\_fp}
\end{itemize}

\cpath{inputs\_test\_2d\_prev\_positions\_picmi.py} 则验证：

\begin{itemize}
\item \cpath{warpx\_save\_previous\_position=True}
\item \cpath{prev\_x/prev\_z} runtime attributes 确实被加进 species
\item \cpath{PushPX()} 在推进前确实保存了旧位置
\end{itemize}

这组输入不检查某条特定的场或轨道物理，而是检查：

\begin{itemize}
\item Python 对 runtime attributes 的增删访问
\item Python 注入粒子接口
\item Python 手动沉积接口
\item Python 到 C++ \cpath{save\_previous\_position} 运行时属性链
\end{itemize}

源码数据流也应分三层理解：PICMI \cpath{sim.particles} 返回 pybind 暴露的 \cpath{WarpX::GetPartContainer()}；species 操作最终调用 \cpath{WarpXParticleContainer.cpp} 的 \cpath{add\_n\_particles(...)}、\cpath{deposit\_current(...)} 或 \cpath{deposit\_charge(...)}；\cpath{beforestep}/\cpath{afterstep} 回调先进入 Python \cpath{CallbackFunctions}，再由 \cpath{ExecutePythonCallback(name)} 调用。因而这组输入检查的是 PICMI 外观、pybind 粒子对象和 callback bridge 的联合行为，而不是某个孤立 Python helper。

这里有一个明确的覆盖缺口：\cpath{test\_2d\_particle\_attr\_access\_unique\_picmi} 虽传入 \cpath{--unique}，输入脚本的 \cpath{add\_particles(...)} 仍硬编码 \cpath{unique\_particles=True}，没有消费 \cpath{args.unique}。因此它不能证明 unique/non-unique 两种注入语义都已经比较过。

同一条接口线还区分两类 restart 读取。\cpath{id\_cpu\_read} 输入遍历 \code{pti["idcpu"]}，用 \cpath{unpack\_ids/unpack\_cpus} 检查累计和为 \cpath{5050/0}，所以它检查的是粒子 id/cpu 的解包读取。runtime-components 输入把 \code{add_real_comp("newPid")}、callback 注入、\code{get_real_comp_index("newPid")} 和 \cpath{picmi.Checkpoint(...)} 放在一起，说明动态 component 可以进入 checkpoint 前端；其 restart 目标虽然存在，但 analysis/checksum 都是 \cpath{OFF}，因此不能据此断言 restart 后 runtime attributes 已获完整验证。

\subsubsection{单精度：有分析程序不等于已有活跃回归}

\cpath{particle\_fields\_diags} 的 single-precision 变体要单独理解：

\begin{itemize}
\item \cpath{analysis\_particle\_diags\_single.py} 已经存在
\item 它复用同一实现，只把容差放宽到 \cpath{5e-3}
\item 但 \cpath{CMakeLists.txt} 里整条 \cpath{test\_3d\_particle\_fields\_diags\_single\_precision} 仍被 \codeesc{\# FIXME} 注释掉
\end{itemize}

因此可以支持的结论是：

\begin{itemize}
\item single-precision particle-field reductions 的 analysis 预案已经随源码提供
\item 但它还不是活跃 regression
\end{itemize}

把这三组放在一起，读者就能避免三种常见误读：粒子从主容器消失不等于它可被 Python 取回；Python 能写入属性不等于 restart 语义已经完整验证；源码提供单精度分析程序不等于该变体正在被自动回归覆盖。

\subsection{边界上的粒子：内建条件、Python 回调与可验证的物理后果}

边界不是一个统一的“反射开关”。在 \cpath{particle\_boundary\_scrape} 和 \cpath{particle\_data\_python} 之外，读者至少要区分三件事：内建边界怎样更新粒子，scraped buffer 能交给用户多少撞击信息，以及 Python 参数前端是否仍把同一物理模型送进 C++ 主链。下面四组测试分别覆盖这些问题：

\begin{itemize}
\item \cpath{particle\_boundary\_interaction}
\item \cpath{particle\_boundary\_process}
\item \cpath{particle\_thermal\_boundary}
\item \cpath{plasma\_lens\_python}
\end{itemize}

它们不只看最终轨道，而是分别通过 Python callback、buffer、parser 或 reduced diagnostics，检查边界语义怎样暴露给用户。

\subsubsection{自己写边界物理：scraped buffer 记录撞击事件，不提供反射后的粒子}

\cpath{particle\_boundary\_interaction} 的关键点首先不是 analysis，而是输入脚本本身的结构。它不是直接依赖 WarpX 内建的 EB 反射，而是：

\begin{enumerate}
\item 把电子打到一个球形 embedded boundary 上
\item 打开 \cpath{warpx\_save\_particles\_at\_eb=1}
\item 在 \cpath{afterstep} callback 里用 \cpath{ParticleBoundaryBufferWrapper.get\_particle\_scraped\_this\_step(...)} 取出：
\begin{itemize}
\item \cpath{deltaTimeScraped}
\item \cpath{r/theta/z}
\item \cpath{ux/uy/uz}
\item \cpath{nx/ny/nz}
\end{itemize}
\item 在 Python 里手工做镜面反射
\item 再按 \code{dt - delta_t} 把粒子推进到这一步的末尾并重新 \cpath{add\_particles(...)}
\end{enumerate}

因此它验证的是：scraped buffer 是否给出了足够的几何和时间信息，让用户实现此例的边界相互作用模型。它记录的是接触时刻、接触位置和局部法向，而不是可以直接继续推进的“反射后粒子”；用户回调仍须决定反射、二次发射或吸收后的动量，并补完 \cpath{dt-deltaTimeScraped}。后面的 \cpath{analysis.py} 用解析几何反射轨道比较最终 \cpath{x/z}，所以能检验该再注入模型的几何正确性，不能证明任意用户回调或任意材料模型正确。

\subsubsection{内建 domain boundary：分别检查粒子数、动量和位置}

在这些更复杂的边界 regression 之外，源码中还有一个更“教科书式”的最小强基准：\cpath{Examples/Tests/boundaries/}。它不碰 embedded boundary、不碰 callback，也不依赖 reduced diagnostics，而是把三类 domain particle boundary 直接拆开测试：

\begin{itemize}
\item \cpath{x} 方向 reflecting
\item \cpath{y} 方向 absorbing
\item \cpath{z} 方向 periodic
\end{itemize}

输入里三类 species 都用 \cpath{MultipleParticles} 直接给定初始位置和动量；analysis 则显式检查：

\begin{itemize}
\item absorbing species 最终只剩 1 个粒子
\item reflecting species 的速度严格翻号
\item periodic species 的速度保持不变
\item 两类保留粒子的位置都与解析反射/wrap-around 位置逐点一致
\end{itemize}

因此 \cpath{particle\_boundaries} 验证的不是应用级“边界大致工作”，而是 \cpath{ParticleBoundaries\_K.H} 最核心的三种 domain particle boundary 语义本体。

\cpath{particle\_boundary\_process} 又不同，它分成两条不同强度的检查。

第一条是 \cpath{test\_2d\_particle\_reflection\_picmi}。虽然 \code{analysis = OFF}，但输入脚本本身做了直接自检：

\begin{itemize}
\item \code{warpx_reflection_model_zhi = "0.5"}
\item 打开 \cpath{save\_particles\_at\_zhi/zlo}
\item 跑完后直接检查：
\begin{itemize}
\item \cpath{z\_hi} buffer 中粒子数是 63
\item \cpath{z\_lo} buffer 中粒子数是 67
\item \cpath{z\_hi} 的 \cpath{stepScraped} 全都等于 4
\item \cpath{z\_lo} 的 \cpath{stepScraped} 全都等于 8
\end{itemize}
\end{itemize}

这组输入内自检测的是：

\begin{itemize}
\item absorbing boundary 上的随机反射模型 parser
\item 上下边界 scraped buffer 的分流
\item \cpath{stepScraped} 时间戳语义
\end{itemize}

第二条是 \cpath{test\_3d\_particle\_absorption}。这里 analysis 很简单，只检查：

\begin{itemize}
\item 第 40 步仍有 612 个电子
\item 第 60 步电子数变成 0
\end{itemize}

所以它更接近“EB 吸收后主 species 中粒子确实消失”的强吸收断言。

\cpath{particle\_thermal\_boundary} 也必须和前面已经整理过的 \cpath{ParticleThermalizer} 区分开。这里输入打开的是：

\begin{itemize}
\item \code{boundary.particle_lo = thermal thermal}
\item \code{boundary.particle_hi = thermal thermal}
\item \code{boundary.<species>.u_th = ...}
\end{itemize}

也就是 domain particle boundary 本身就是 thermal boundary。对应 \cpath{ParticleBoundaries\_K.H} 的语义是：先像反射边界一样处理位置，再对法向/切向动量做热化抽样。analysis 并不看单粒子散射角，而是读取 reduced diagnostics 的：

\begin{itemize}
\item \cpath{FieldEnergy}
\item \cpath{ParticleEnergy}
\end{itemize}

并要求：

\begin{itemize}
\item 场能不能无界增长
\item 粒子总能量相对初值偏离不超过 2\%
\end{itemize}

所以它检验的是 thermal particle boundary 在长时间粒子出入边界时的总量稳定性；它并不是单粒子散射角分布的充分验证。

\subsubsection{Python 参数前端：同一物理断言，增加的是参数传递覆盖}

最后，\cpath{plasma\_lens\_python} 复用和 native/PICMI plasma-lens 相同的 \cpath{analysis.py}，所以物理断言没有变化，仍然是两颗测试电子穿过 lens 序列后的：

\begin{itemize}
\item 最终横向位置
\item 最终横向动量
\end{itemize}

与解析模型一致。

但它新增覆盖的不是物理，而是 front-end：输入不再走 native 文件或 PICMI，而是直接用 \cpath{pywarpx} 参数对象设置：

\begin{itemize}
\item \code{particles.E_ext_particle_init_style = "repeated_plasma_lens"}
\item \code{particles.B_ext_particle_init_style = "repeated_plasma_lens"}
\item \code{particles.repeated_plasma_lens_* = ...}
\item 最后 \code{warpx.init(); warpx.step(...)}
\end{itemize}

它新增覆盖的是纯 Python 参数前端到 \cpath{MultiParticleContainer} / \cpath{GetExternalEBField} repeated-plasma-lens 主链，而不是新的 lens 物理。阅读这些例子时，应按问题选择观察量：写 Python 表面物理就检查撞击几何、法向和剩余时间；确认内建边界就分别检查粒子数、动量和解析位置；确认 Python front-end 就使用与 native 例相同的物理观察量。三种证据不能互相替代。

同一条 Python scraped-buffer 物理链还有一个更直接的边界物理 regression：\cpath{secondary\_ion\_emission}。它不是只拿 buffer 做统计，而是在 \cpath{afterstep} callback 里直接用

\begin{itemize}
\item \cpath{r/theta/z}
\item \cpath{ux/uy/uz}
\item \cpath{nx/ny/nz}
\item \cpath{deltaTimeScraped}
\end{itemize}

为撞击球形 EB 的离子生成次级电子。analysis 再要求：

\begin{enumerate}
\item 固定随机种子下最终恰好产生 2 个电子；
\item 电子反向传播到撞击时刻后，应落在解析球面撞击点附近
\end{enumerate}

所以这组例子说明：\cpath{ParticleBoundaryBufferWrapper} 的几何和时间元数据足以支撑该 callback 驱动的二次发射模型，而不只是后处理统计；它不自动给出通用的表面材料模型。

\subsection{接触记录、PML 粒子与 AMR：观察对象决定验证强度}

embedded boundary、PML 和 mesh refinement 都会改变粒子所在的数值环境，但它们的正确性不能用同一个观察量判断。下面四类例子分别检查接触几何、残余场、组合稳定性和解析场一致性：

\begin{itemize}
\item \cpath{point\_of\_contact\_eb}
\item \cpath{particles\_in\_pml}
\item \cpath{subcycling\_mr}
\item \code{Langmuir multi_mr}
\end{itemize}

输入、analysis 和 writer 路径表明，它们回答的是四个不同的问题。

\subsubsection{接触记录：确认写出的事件几何，而不是只确认粒子消失}

\cpath{point\_of\_contact\_eb} 的关键是它打开了两套 diagnostics：

\begin{itemize}
\item \code{diag1 = Full}
\item \code{diag2 = BoundaryScraping}
\end{itemize}

analysis 根本不看主 species 的最终位置，而是直接读：

% <-- fenced code (text)
\begin{codeblock}
diags/diag2/particles_at_eb/
\end{codeblock}

也就是 \cpath{BoundaryScraping} 写出的 scraped-particle openPMD 输出。它比较的是：

\begin{itemize}
\item \cpath{stepScraped}
\item \cpath{deltaTimeScraped}
\item \cpath{x/y/z}
\item \cpath{nx/ny/nz}
\end{itemize}

与解析接触点、接触时刻和表面法向是否一致。因此它验证的是：

\begin{itemize}
\item EB 接触事件是否被正确记录到 \cpath{particles\_at\_eb}
\item \cpath{BoundaryScraping} 输出里的几何量和时间量是否正确
\end{itemize}

这和前面的 \cpath{particle\_boundary\_scrape} 有本质区别：后者更侧重“粒子被删掉并进了 buffer”，而 \cpath{point\_of\_contact\_eb} 检查记录下来的接触几何和时间是否正确。

\subsubsection{粒子进入 PML：用残余场检验清理是否有效}

\cpath{particles\_in\_pml} 则不是看粒子轨道，而是看粒子离域后留下的场。analysis 的核心只有一句：

% <-- fenced code (python)
\begin{codeblock}
assert max_Efield < tolerance_abs
\end{codeblock}

它先取最后一步的 \cpath{Ex/Ey/Ez}，再要求域内残余电场足够小。对应物理含义是：

\begin{itemize}
\item 如果 PML 只吸场不吸粒子，离开的带电粒子会留下伪电荷和伪场
\item 打开 \code{warpx.pml_has_particles = 1} 后，这个 spurious field 必须被压到足够低
\end{itemize}

2D/3D 与 MR 版本共享同一 analysis，只是 tolerance 按：

\begin{itemize}
\item 维度
\item \cpath{max\_level}
\end{itemize}

分开设置。因此它检验的是 particle-aware PML 的 residual-field cleanup，而不是单纯的 PML 场反射率；比较不同维度或不同 AMR level 时，不能忽略容差本身的差异。

\subsubsection{AMR 与 subcycling：checksum 基线和解析基准的证据强度不同}

\cpath{subcycling\_mr} 又更弱一层。它当前没有独立 analysis，只有 checksum。输入里同时打开了：

\begin{itemize}
\item \code{warpx.do_subcycling = 1}
\item \code{amr.max_level = 1}
\item moving window
\item driver / beam / plasma continuous injection
\item \code{particles.deposit_on_main_grid = plasma_e plasma_p}
\item \code{n_current_deposition_buffer = 0}
\item \code{n_field_gather_buffer = 0}
\end{itemize}

所以它只能建立 \code{AMR + subcycling + moving window + continuous injection + deposit_on_main_grid} 这一组合的 checksum 基线。由于没有独立观察量，它不能分辨 injection、gather、deposit 或 coarse/fine 同步的任何一个步骤，也不能写成独立的 refined-injection analysis。

最后，\code{Langmuir multi_mr} 与 \cpath{subcycling\_mr} 正好相反。它虽然也在测 MR，但并不只是 checksum，而是继续复用 \cpath{analysis\_2d.py} 的强验证：

\begin{enumerate}
\item 取 \cpath{Ex/Ez}
\item 与解析 Langmuir-wave 场解比较
\item 在满足条件时，再检查 \cpath{divE} 与 \cpath{rho/eps0} 的相对误差
\end{enumerate}

而不同 MR 变体测的是不同 refinement 组合：

\begin{itemize}
\item \cpath{mr}
\begin{itemize}
\item \code{amr.max_level = 1}
\item \code{amr.ref_ratio = 4}
\end{itemize}
\item \cpath{mr\_anisotropic}
\begin{itemize}
\item \code{amr.ref_ratio_vect = 4 2}
\end{itemize}
\item \cpath{mr\_maxlevel2}
\begin{itemize}
\item \code{amr.max_level = 2}
\end{itemize}
\item \cpath{mr\_momentum\_conserving}
\begin{itemize}
\item gather scheme 改成 momentum-conserving
\end{itemize}
\item \cpath{mr\_psatd}
\begin{itemize}
\item solver 改成 PSATD
\end{itemize}
\end{itemize}

因此它可支持“这些已覆盖的 MR 组合没有破坏该解析 Langmuir 粒子-场基准”；不能外推为任意 AMR 配置、任意注入模型或所有 solver 的证明。

把四类例子放在同一张心理地图中最有用：接触 writer 看几何与时间，particle-aware PML 看残余场，subcycling 例提供组合 checksum，而 Langmuir multi-MR 保留解析场和电荷关系的强基准。观察对象不同，能得出的结论也必须不同。

实际阅读或设计自己的测试时，可以按下列顺序选择观察量：

\begin{enumerate}
\item 若问题是“粒子何时、何处触及表面”，优先读取 \cpath{stepScraped}、\cpath{deltaTimeScraped}、位置和法向；只数剩余粒子无法判断接触几何。
\item 若问题是“离开物理域的带电粒子是否留下数值污染”，比较 PML 后时刻的残余 \cpath{Ex/Ey/Ez}，不要把它与入射波的反射率混为一谈。
\item 若问题是“一个复杂 AMR 输入能否保持可重复输出”，checksum 是必要的基线；若要判断场和电荷关系是否仍正确，则必须再选择有解析解或守恒量的基准，如 Langmuir multi-MR。
\item 若某个例子只提供 checksum，读者应把它当作该配置组合的漂移报警器，而不是将其当成每一条粒子、网格和同步路径都已经独立证明。
\end{enumerate}

这套区分也解释了为什么一个通过的测试集合仍可能留下开放边界：它们的输入可以相似，但被比较的对象、容差和可外推的范围并不相同。

\subsection{Embedded boundary：从几何、吸收和解析场选择验证量}

要验证 embedded boundary，先要确定你想验证的是几何表示、场的边界条件、粒子吸收，还是 Python 对几何数组的访问。下列例子涉及这些不同问题：

\begin{itemize}
\item \cpath{particle\_absorbing\_boundary/plot\_2d.py}
\item \cpath{particle\_absorbing\_boundary/plot\_phase.py}
\item \cpath{embedded\_boundary\_cube}
\item \cpath{embedded\_boundary\_rotated\_cube}
\item \cpath{embedded\_boundary\_diffraction}
\item \cpath{embedded\_boundary\_em\_particle\_absorption}
\item \cpath{embedded\_boundary\_python\_api}
\item \cpath{electrostatic\_sphere\_eb}
\item \cpath{scraping}
\end{itemize}

先分清工具和断言：\cpath{plot\_2d.py} 与 \cpath{plot\_phase.py} 都不是 regression analysis。

它们只是：

\begin{itemize}
\item 画 2D full diagnostics 的 \cpath{Ez} slice
\item 画 \cpath{PhaseSpaceElectrons} reduced diagnostic 的相图
\end{itemize}

它们只是 visualization helper，不是物理断言脚本；图像适合帮助读者定位异常，却不能代替定量误差或守恒检查。

\subsubsection{PEC cavity：用解析本征模检验几何与场边界}

第一类是 cavity 模态解析对照：

\begin{itemize}
\item \cpath{embedded\_boundary\_cube}
\item \cpath{embedded\_boundary\_rotated\_cube}
\end{itemize}

\cpath{analysis\_fields.py}、\cpath{analysis\_fields\_2d.py} 和 \cpath{analysis\_fields\_3d.py} 都是显式构造 PEC cavity 的解析本征模，再比较 \cpath{By}、\cpath{Bz} 或 \cpath{Ey/c} 的相对 \cpath{L2} 误差。区别只是：

\begin{itemize}
\item \cpath{cube}：轴对齐 cavity
\item \cpath{rotated\_cube}：旋转后的 cavity，analysis 里要先把坐标和场分量反旋回解析坐标系
\item \cpath{cube\_macroscopic}：同一模态，但频率按介质 \cpath{epsilon\_r} 修正
\end{itemize}

这两组例子回答的是“给定的 EB 几何和 PEC 场边界能否维持指定本征模”。旋转例还把坐标和场分量转回解析坐标系，因而同时约束几何方向和矢量分量的处理；它们不直接验证带电粒子撞击或任意复杂几何。

\subsubsection{衍射与 Python API：一个测物理图样，一个测几何数组}

\begin{itemize}
\item \cpath{embedded\_boundary\_diffraction}
\item \cpath{embedded\_boundary\_python\_api}
\end{itemize}

\cpath{embedded\_boundary\_diffraction/analysis\_fields.py} 读取 RZ 下的 \cpath{Ex}，提取衍射图样第一极小值半径，再与 Airy pattern 的
\[
\theta \sim 1.22 \lambda / d
\]
预测比较。因此它测的是 EB 产生的衍射图样是否在 Airy 第一极小值上与解析预测相符，而不是对每个网格单元的几何误差证明。

\cpath{embedded\_boundary\_python\_api} 更特殊。CMake 里虽然 \code{analysis = OFF}，但 PICMI 输入脚本本身在运行时就会读取：

\begin{itemize}
\item \cpath{edge\_lengths}
\item \cpath{face\_areas}
\end{itemize}

并在三个中间切片上重建 cavity 的 perimeter 和 area，再和解析几何值比较。它虽没有独立 \cpath{analysis.py}，却不是单纯 checksum：输入本身验证了 PICMI wrapper 返回的 edge-length 与 face-area 几何量。这个检查不自动覆盖任意 EB 形状、所有 AMR level 或后续场演化。

\subsubsection{吸收与 scraping：分别观察伪电荷和粒子记账}

\begin{itemize}
\item \cpath{embedded\_boundary\_em\_particle\_absorption}
\item \cpath{scraping}
\end{itemize}

\cpath{embedded\_boundary\_em\_particle\_absorption/analysis.py} 做的不是看粒子是否消失，而是把 \cpath{divE} 做时间平均，去掉沿 EB 传播的真实波动分量后，检查是否还残留静态伪电荷。因此它检验的是 EM 粒子吸收不会在该测试条件下积累非物理电荷，而不是一般的 \cpath{divE-rho} 全时空闭合证明。

而 \cpath{scraping/analysis\_rz.py} 与 \cpath{analysis\_rz\_filter.py} 测的是另一条 writer 合同：

\begin{itemize}
\item 最终剩余粒子数是否正确
\item \code{remaining + scraped = initial} 是否逐步成立
\item scraped buffer 中的 \cpath{id} 是否和初始全集闭合
\item 打开 \cpath{plot\_filter\_function} 后，是否真的只记录 \code{z > 0} 半域的 scraped particles
\end{itemize}

因此 \cpath{scraping} 关注的是 BoundaryScraping 的粒子记账和 \cpath{plot\_filter\_function} 的选择语义。它应与前一例的场量检查分开阅读：前者能发现 particle identity 或输出选择错误，后者能发现吸收后的静态场污染。

\subsubsection{静电球：用解析势与电荷把三维、RZ 和 AMR 分层比较}

\cpath{electrostatic\_sphere\_eb} 至少分成三层：

\begin{enumerate}
\item 3D \cpath{analysis.py}
\begin{itemize}
\item reduced diagnostic \cpath{eb\_charge.txt}
\item 理论球导体电荷
\end{itemize}
\end{enumerate}
\[
     q = 4 \pi \epsilon_0 \phi_0 R
\]
\begin{itemize}
\item \cpath{eb\_covered} 场是否满足内 \cpath{1} 外 \cpath{0}
\end{itemize}
\begin{enumerate}
\item RZ \cpath{analysis\_rz.py}
\begin{itemize}
\item 比较解析
\end{itemize}
\end{enumerate}
\[
     \phi(r)=A+B\log r,\qquad E_r(r)=-B/r
\]
\begin{enumerate}
\item RZ \cpath{analysis\_rz\_mr.py}
\begin{itemize}
\item 把同一 $\phi/E_r$ 对照扩展到每个 refinement level
\end{itemize}
\end{enumerate}

只有 \cpath{inputs\_test\_3d\_electrostatic\_sphere\_eb\_mixed\_bc} 没有独立 analysis，因此它只提供 mixed-BC 配置的 checksum 基线。其余例子分别把导体总电荷、覆盖区域、RZ 解析势/径向场和各 refinement level 的误差作为观察量；不能把通过的 checksum 当作解析 $\phi/E_r$ 比较，也不能把一个对称球的结果外推为任意电极几何。

综上，embedded-boundary 验证至少应同时问四个问题：几何是否表示正确，场是否满足已知解析解，吸收是否留下数值电荷，Python 或 writer 是否导出了预期的几何/粒子数据。不同问题需要不同的输出量，任何一个单独通过都不等于整个 EB 物理闭合。

\section{QED：先分清 source、product 与产生机制，再谈参数}

从 \cpath{Particles/} 入口再往下看，WarpX 当前的 QED 主链至少分成三种完全不同的事件类型：

\begin{enumerate}
\item Quantum Synchrotron：lepton source 产生 photon product
\item Breit-Wheeler：photon source 产生 electron/positron products
\item Schwinger：不以 source species 为起点，而是直接由场在网格上创建对
\end{enumerate}

\subsection{事件开始前：谁是 source，谁保存统计状态，谁接收 product}

\cpath{PhysicalParticleContainer.cpp} 在构造期先决定：

% <-- fenced code (cpp)
\begin{codeblock}
pp_species_name.query("do_qed_quantum_sync", m_do_qed_quantum_sync);
if (m_do_qed_quantum_sync) {
    AddRealComp("opticalDepthQSR");
}

pp_species_name.query("do_qed_breit_wheeler", m_do_qed_breit_wheeler);
if (m_do_qed_breit_wheeler) {
    AddRealComp("opticalDepthBW");
}
\end{codeblock}

这不是实现细节，而是 QED 与 field ionization 一样，先把“事件统计状态”写进 species 的 runtime attribute 系统。读者随后追踪每一次事件时，应先找到：

\begin{itemize}
\item \cpath{opticalDepthQSR}
\item \cpath{opticalDepthBW}
\end{itemize}

这两个持久属性。

与之配套，\cpath{PhotonParticleContainer.cpp} 又明确禁止 photon species 再开 quantum synchrotron：

% <-- fenced code (cpp)
\begin{codeblock}
WARPX_ALWAYS_ASSERT_WITH_MESSAGE(
    test_quantum_sync == 0,
    "ERROR: do_qed_quantum_sync can't be enabled for photon particles!");
\end{codeblock}

因此，WarpX 在 species 层已经把“谁提供待转化粒子、谁保存 optical depth、谁接收新粒子”分开；只有先分清这些角色，才能正确解读后续的粒子数、动量与权重变化。

\subsection{初始化只准备采样条件，事件要在后续时间步发生}

\cpath{MultiParticleContainer::InitMultiPhysicsModules()} 在 \cpath{InitData()}/\cpath{PostRestart()} 前先做：

% <-- fenced code (cpp)
\begin{codeblock}
mapSpeciesProduct();
CheckQEDProductSpecies();
InitQED();
\end{codeblock}

这里三步分别对应：

\begin{enumerate}
\item 把三个 product-species 参数从字符串映射成容器索引：
\begin{itemize}
\item \cpath{qed\_quantum\_sync\_phot\_product\_species}
\item \cpath{qed\_breit\_wheeler\_ele\_product\_species}
\item \cpath{qed\_breit\_wheeler\_pos\_product\_species}
\end{itemize}
\item 检查 product species 类型是否正确；
\item 创建 \cpath{QuantumSynchrotronEngine} / \cpath{BreitWheelerEngine} 并按 \cpath{qed\_qs.*}、\cpath{qed\_bw.*} 初始化 lookup tables。
\end{enumerate}

因此 \cpath{InitQED()} 不是“提前跑一遍 QED”，而是在模拟开始或 restart 后把下列采样条件固定：

\begin{itemize}
\item 谁参与 quantum synchrotron
\item 谁参与 Breit-Wheeler
\item 表从 builtin/load/generate 哪条路径来
\end{itemize}

这一区分决定了排错顺序：初始化失败时先检查 product species 与 lookup table；粒子数或能谱异常时才追踪后续 push 与 event pass。

\subsection{时间位置：在 field ionization 后、用户 injection 前处理 QED}

顶层主循环 \cpath{WarpXEvolve.cpp} 里，多物理顺序非常直接：

% <-- fenced code (cpp)
\begin{codeblock}
doFieldIonization();

#ifdef WARPX_QED
doQEDEvents();
mypc->doQEDSchwinger();
#endif

ExecutePythonCallback("particleinjection");
OneStep(cur_time, dt[0], step);
\end{codeblock}

这说明 QED 事件发生在：

\begin{itemize}
\item field ionization 之后
\item 用户 \cpath{particleinjection} callback 之前
\item 正常 \cpath{OneStep()} 之前
\end{itemize}

这说明 QED 不像 collisions 那样嵌在 split-momentum push 的组织里，而是更早地消费当下的 \cpath{Efield\_aux/Bfield\_aux}。因此，若用户 callback 在 \cpath{particleinjection} 中新增粒子，这些粒子不会在同一轮的 QED event pass 中被处理。

\subsection{两条 source-to-product 事件链：转化时同时更新 source 与创建 product}

\cpath{doQedQuantumSync()} 的骨架是：

% <-- fenced code (cpp)
\begin{codeblock}
const auto Filter   = phys_pc_ptr->getPhotonEmissionFilterFunc();
const auto CopyPhot = copy_factory_phot.getSmartCopy();

auto Transform = PhotonEmissionTransformFunc(
      m_shr_p_qs_engine->build_optical_depth_functor(),
      pc_source->GetRealCompIndex("opticalDepthQSR") - pc_source->NArrayReal,
      m_shr_p_qs_engine->build_phot_em_functor(),
      pti, lev, Ex.nGrowVect(), ...);

filterCopyTransformParticles<1>(
    *pc_product_phot, dst_tile, src_tile, np_dst,
    Filter, CopyPhot, Transform);
\end{codeblock}

\cpath{doQedBreitWheeler()} 则是双 product 版本：

% <-- fenced code (cpp)
\begin{codeblock}
const auto Filter  = phys_pc_ptr->getPairGenerationFilterFunc();
const auto pair_gen_functor = m_shr_p_bw_engine->build_pair_functor();

auto Transform = PairGenerationTransformFunc(pair_gen_functor,
                                             pti, lev, Ex.nGrowVect(), ...);

filterCopyTransformParticles<1>(
    *pc_product_ele, *pc_product_pos,
    dst_ele_tile, dst_pos_tile, src_tile,
    np_dst_ele, np_dst_pos,
    Filter, CopyEle, CopyPos, Transform);
\end{codeblock}

这两条链有一个关键共同点：QED 不是只在 product species 上增加一批新粒子，而是同时做三件事：

\begin{enumerate}
\item 用 source 粒子的 optical depth 判断是否触发事件
\item 在 \cpath{Transform} 里读取主网格场和 external particle fields
\item 同时更新 source 状态并创建 product particles
\end{enumerate}

这与前面 field ionization 的 \cpath{filterCopyTransformParticles} 思路相近，但它依赖 QED engine 表和 optical-depth 统计变量。验证时因此不能只计数 product，还要同时检查 source 的剩余状态与能量/动量变化。

\subsection{Schwinger：由网格场直接创建粒子对，不存在 source species}

\cpath{doQEDSchwinger()} 不再从某个 source species 复制，而是直接在网格上用：

\begin{itemize}
\item \cpath{Efield\_aux/Bfield\_aux}
\item Schwinger 激活区域
\item \cpath{filterCreateTransformFromFAB<1>(...)}
\end{itemize}

生成 \cpath{ele\_schwinger} 和 \cpath{pos\_schwinger}。而且它目前硬性要求：

\begin{itemize}
\item collocated grid 或 momentum-conserving gather
\item 无 mesh refinement
\item 非 RZ
\item 非 1D
\end{itemize}

所以不能把 Schwinger 和前两条 source-to-product 路径混写成同一类机制：它的观察起点是激活区域内的场和 pair 产额，而不是某一 source species 的 optical depth。

\subsection{三条机制分别该看什么：数目之外还要看权重、方向和守恒}

\begin{itemize}
\item \cpath{analysis\_quantum\_sync.py}：
\begin{itemize}
\item 检查 photon 数、权重、发射方向、能谱、以及 source/product optical depth 分布；
\item 对应 Quantum Synchrotron 的主合同。
\end{itemize}
\item \cpath{analysis\_breit\_wheeler\_core.py}：
\begin{itemize}
\item 检查 pairs 数、残余 photon 动量、单事件能量守恒、pair 能谱和 optical depth；
\item 对应 Breit-Wheeler 的 product-species 合同。
\end{itemize}
\item \cpath{analysis\_schwinger.py}：
\begin{itemize}
\item 检查理论率对应的 pair 数窗口，以及电子/正电子权重数组一致性；
\item 对应 Schwinger 的强场真空对产生合同。
\end{itemize}
\end{itemize}

因此“QED 测试通过”不是单一结论：Quantum Synchrotron 需要同时看 photon 发射与 lepton 的统计状态，Breit-Wheeler 需要比较残余 photon、pairs 与单事件守恒，Schwinger 则以理论率窗口和电子/正电子权重配对为主。三条路径必须分别验证，不能用其中一条替代另外两条。

\subsection{事件何时发生：先演化 optical depth，再决定是否创建 product}

把入口层再往下读到 \cpath{QEDPhotonEmission.H}、\cpath{QEDPairGeneration.H} 和两个 engine wrapper，会发现 QED 事件触发并不是“在 event pass 里直接抽一次随机数”。

\cpath{PhotonEmissionFilterFunc} 和 \cpath{PairGenerationFilterFunc} 都只做一件事：

% <-- fenced code (cpp)
\begin{codeblock}
return (opt_depth < 0.0_rt);
\end{codeblock}

这表示 event pass 不负责从零开始抽样，而是消费先前 push 已更新的状态。真正的统计演化发生在更早的 push 阶段：

\begin{itemize}
\item \cpath{PhysicalParticleContainer::PushPX()} 里先用 \cpath{QuantumSynchrotronEvolveOpticalDepth} 推进 \cpath{opticalDepthQSR}
\item \cpath{PhotonParticleContainer::PushPX()} 里先用 \cpath{BreitWheelerEvolveOpticalDepth} 推进 \cpath{opticalDepthBW}
\end{itemize}

只有当 optical depth 在 push 中被推进到负值，后面的：

\begin{itemize}
\item \cpath{doQedQuantumSync()}
\item \cpath{doQedBreitWheeler()}
\end{itemize}

才会在 event pass 里通过 \cpath{filterCopyTransformParticles} 触发事件。调试“没有产生 photon/pair”时，应先检查 \cpath{chi}、时间步和 optical depth 是否能跨过零，而不是只检查 product species 是否存在。

\subsection{分工边界：WarpX 提供粒子与场，PICSAR-QED 提供采样内核}

\cpath{QuantumSyncEngineWrapper.H} 和 \cpath{BreitWheelerEngineWrapper.H} 不是又实现了一遍 QED 理论，而是把 PICSAR-QED core 包成 WarpX 可在 GPU kernel 中调用的 functor：

\begin{itemize}
\item \cpath{QuantumSynchrotronEvolveOpticalDepth} 最终调用 \cpath{pxr\_qs::evolve\_optical\_depth}
\item \cpath{QuantumSynchrotronPhotonEmission} 最终调用 \cpath{pxr\_qs::generate\_photon\_update\_momentum}
\item \cpath{BreitWheelerEvolveOpticalDepth} 最终调用 \cpath{pxr\_bw::evolve\_optical\_depth}
\item \cpath{BreitWheelerGeneratePairs} 最终调用 \cpath{pxr\_bw::generate\_breit\_wheeler\_pairs}
\end{itemize}

WarpX 在这一层真正负责的是：

\begin{enumerate}
\item gather 主网格与 external particle fields
\item 提供 source/product species 的 SoA 指针
\item 存取 \cpath{opticalDepthQSR/BW}
\item 组织 \cpath{filterCopyTransformParticles}
\end{enumerate}

这条分工线给出了排错边界：gather、species 属性、tile 调度和 product 容器属于 WarpX；给定 \cpath{chi} 后的概率演化与动量采样属于 PICSAR-QED。不要把 table 或采样偏差误诊为 WarpX 的粒子容器问题，也不要把错误的 field gather 归咎于采样器。

\subsection{source 的命运不同：辐射会保留 lepton，成对产生会消耗 photon}

两条 event kernel 都会先 gather \cpath{E/B}，但 source 处理方式并不一样：

\begin{itemize}
\item \cpath{PhotonEmissionTransformFunc} 会原地改写 source lepton 动量，并把 source optical depth 重新抽样初始化；
\item \cpath{PairGenerationTransformFunc} 会生成 electron/positron 两个 product，并把 source photon 直接标记成 invalid。
\end{itemize}

这解释了为什么两类分析需要不同的守恒与统计检查：

\begin{itemize}
\item \cpath{analysis\_quantum\_sync.py} 要重点检查 source/product optical-depth 分布能否继续保持指数；
\item \cpath{analysis\_breit\_wheeler\_core.py} 要重点检查 residual photons、丢失 photon 数和新 pairs 数之间的对应关系。
\end{itemize}

\subsection{lookup table 是采样器的一部分，不是可有可无的附属文件}

如果继续往 \cpath{QEDInternals/QuantumSyncEngineWrapper.cpp} 和 \cpath{BreitWheelerEngineWrapper.cpp} 读，会发现 WarpX 当前的 QED wrapper 内部都不是只持有一张表，而是：

\begin{itemize}
\item 一张 \cpath{dndt} 表
\item 一张真正用于事件采样的二维表
\item 一个最小 \cpath{chi} 门槛
\end{itemize}

更具体地说：

\begin{itemize}
\item Quantum Synchrotron 持有 \code{m_dndt_table + m_phot_em_table + m_qs_minimum_chi_part}
\item Breit-Wheeler 持有 \code{m_dndt_table + m_pair_prod_table + m_bw_minimum_chi_phot}
\end{itemize}

因此前面看到的：

\begin{itemize}
\item \cpath{build\_optical\_depth\_functor()}
\item \cpath{build\_phot\_em\_functor()}
\item \cpath{build\_pair\_functor()}
\end{itemize}

都会先要求 \code{m_lookup_tables_initialized == true}。这说明 QED kernel 不是“有 wrapper 就能跑”，而是必须先完成表的生命周期；出现初始化或运行时 table 错误时，应先确认模式、输入文件和最小 \cpath{chi}，再解释粒子统计结果。

\cpath{InitQED()} 里 \cpath{qed\_qs.lookup\_table\_mode} 和 \cpath{qed\_bw.lookup\_table\_mode} 只允许三种模式：

\begin{itemize}
\item \cpath{builtin}
\item \cpath{load}
\item \cpath{generate}
\end{itemize}

其中：

\begin{itemize}
\item \cpath{builtin} 直接使用 wrapper \cpath{.cpp} 中硬编码的低分辨率测试表
\item \cpath{load} 从外部二进制文件读 raw bytes，再反序列化成两张子表
\item \cpath{generate} 运行时现生成表，但最后仍然导出成同一种 raw binary 格式，并通过 MPI 广播给所有 rank
\end{itemize}

所以 \cpath{generate} 和 \cpath{load} 的真正共同点是：后续 kernel 消费的不是“生成态”或“加载态”的 C++ 对象，而是同一种序列化结果重新装配出来的 wrapper。

这一点在 \cpath{QuantumSyncGenerateTable()} / \cpath{BreitWheelerGenerateTable()} 里最清楚：

\begin{enumerate}
\item 只在 IOProcessor 上按 \cpath{tab\_dndt\_*}、\cpath{tab\_em\_*} 或 \cpath{tab\_pair\_*} 生成两张子表
\item 调 \cpath{export\_lookup\_tables\_data()} 导出 raw bytes
\item 把 bytes 写入 \cpath{save\_table\_in}
\item 再把同一份 bytes 广播给所有 rank
\item 非 IOProcessor 用 \cpath{init\_lookup\_tables\_from\_raw\_data(...)} 重建 wrapper
\end{enumerate}

因此，完整主链应写成：

% <-- fenced code (text)
\begin{codeblock}
runtime attributes / product species
-> InitQED()
-> builtin/load/generate tables
-> build device functors
-> PushPX() 中演化 optical depth
-> event pass 中查表并生成 product particles
\end{codeblock}

所以 examples 中的 \code{lookup_table_mode = builtin / load / generate} 不只是输入模板，而是直接选择 QED 采样器的可执行上游；同一物理案例切换模式时，仍须核对表分辨率和来源是否适合目标精度。

\subsection{不要被共同的 photon 名称误导：强场 QED 与碰撞 QED 是两棵树}

继续沿源码往下追，会碰到一组很容易被误写进同一章法的名字：

\begin{itemize}
\item \cpath{QedChiFunctions}
\item \cpath{do\_qed\_virtual\_photons}
\item \cpath{linear\_breit\_wheeler}
\item \cpath{linear\_compton}
\end{itemize}

它们都带着 QED / photons 的标签，却不共享同一套事件骨架。选模型前，先问过程由局部强场驱动，还是由 cell 内两粒子配对驱动。

\cpath{QedChiFunctions.H} 本身只有两个薄包装：

\begin{itemize}
\item \cpath{QedUtils::chi\_ele\_pos(...)}
\item \cpath{QedUtils::chi\_photon(...)}
\end{itemize}

它们只把 SI 单位下的动量与 \cpath{E/B} 场交给 PICSAR 的 \cpath{chi} 公式。当前源码里，这两个函数主要只服务三类地方：

\begin{enumerate}
\item \cpath{PushSelector.H} 里 RR 与 QED 联动时的 \cpath{chi} 阈值判断
\item \cpath{QuantumSyncEngineWrapper} 的 optical-depth 与 photon-emission 采样
\item \cpath{BreitWheelerEngineWrapper} 的 optical-depth 与 pair-generation 采样
\end{enumerate}

因此 \cpath{QedChiFunctions} 属于强场 QED 树：

% <-- fenced code (text)
\begin{codeblock}
gather E/B
-> compute chi
-> evolve optical depth
-> lookup-table sampling
-> update source and create products
\end{codeblock}

但 virtual photons 走的不是这条路。

\cpath{PhysicalParticleContainer.cpp} 允许 lepton species 打开：

\begin{itemize}
\item \cpath{do\_qed\_virtual\_photons}
\item \cpath{qed\_virtual\_photon\_species\_name}
\item \cpath{qed\_virtual\_photons\_do\_beam\_size\_effect}
\end{itemize}

随后 \cpath{CollisionHandler::doCollisions()} 在真正做任意碰撞之前，会统一先调用：

% <-- fenced code (cpp)
\begin{codeblock}
collision::binarycollision::virtualphotons::GenerateVirtualPhotons(mypc);
\end{codeblock}

这说明 virtual photons 的运行时位置在 collision 调度层，而不是：

\begin{itemize}
\item \cpath{InitQED()}
\item \cpath{doQedEvents()}
\item \cpath{QEDPhotonEmission.H}
\item \cpath{QEDPairGeneration.H}
\end{itemize}

\cpath{GenerateVirtualPhotons()} 的语义是：每个 coarse step 都从 lepton species 重新采样一批辅助 photon species；旧 virtual photons 会被下一步覆盖。3D 下若打开 beam-size effect，还会在垂直于动量的平面内给这些虚光子加上有限半径位移。

因此 virtual photons 不是强场 QED 事件生成、随后继续 push 的 product photons，而是碰撞模块每个 coarse step 重建的辅助 photon 分布。若分析关注其空间分布或谱，应以 virtual-photon 专用输出和 analysis 为准，不能借用 \cpath{opticalDepthBW} 的解释。

接下来的 \cpath{linear\_breit\_wheeler} 与 \cpath{linear\_compton} 也继续留在碰撞树里。

\cpath{CollisionHandler.cpp} 对 \code{type = linear_breit_wheeler} 的分派是：

% <-- fenced code (cpp)
\begin{codeblock}
std::make_unique<
    BinaryCollision<LinearBreitWheelerCollisionFunc, ParticleCreationFunc>
>(...)
\end{codeblock}

这说明它不走前面的：

\begin{itemize}
\item \cpath{doQedBreitWheeler()}
\item \cpath{PairGenerationFilterFunc}
\item \cpath{BreitWheelerEngineWrapper}
\item \cpath{opticalDepthBW}
\end{itemize}

而是走另一套碰撞骨架：

% <-- fenced code (text)
\begin{codeblock}
optional GenerateVirtualPhotons()
-> BinaryCollision pairing in each cell
-> p_mask / reaction weight
-> ParticleCreationFunc
\end{codeblock}

\cpath{LinearBreitWheelerCollisionFunc} 与 \cpath{LinearComptonCollisionFunc} 自己也写得很清楚：它们实现的是 Higginson 风格的 binary-collision 算法，控制参数是：

\begin{itemize}
\item \cpath{event\_multiplier}
\item \cpath{probability\_threshold}
\item \cpath{probability\_target\_value}
\end{itemize}

而不是强场 QED 那套：

\begin{itemize}
\item \cpath{chi\_min}
\item \cpath{lookup\_table\_mode}
\item \cpath{photon\_creation\_energy\_threshold}
\end{itemize}

因此，WarpX 至少有两棵名字都带 QED / photons 的树：

\begin{enumerate}
\item 强场 QED / \cpath{ElementaryProcess}
\begin{itemize}
\item \cpath{QedChiFunctions}
\item optical depth
\item tables
\item \cpath{filterCopyTransformParticles}
\end{itemize}
\item 碰撞 QED / \cpath{BinaryCollision}
\begin{itemize}
\item optional virtual photons
\item cell-local pairing
\item reaction masks
\item \cpath{ParticleCreationFunc}
\end{itemize}
\end{enumerate}

相应地，测试也分成两组不同的证据：

\begin{itemize}
\item \cpath{analysis\_quantum\_sync.py}、\cpath{analysis\_breit\_wheeler\_core.py}、\cpath{analysis\_schwinger.py}
\begin{itemize}
\item 验证强场 QED 主链
\end{itemize}
\item \cpath{analysis\_virtual\_photons.py}、\cpath{analysis\_beamsize\_effect.py}、\cpath{analysis\_many\_photons.py}
\begin{itemize}
\item 验证 virtual-photon 采样与 linear Breit-Wheeler 碰撞分叉
\end{itemize}
\end{itemize}

实际选择模型时，可用一个简短判据：若观察量是强场中的 \cpath{chi}、optical depth、emission spectrum 或 source photon/lepton 的转化，沿 \cpath{ElementaryProcess} 树阅读；若观察量是两个粒子或虚光子的 cell-local 配对、reaction weight 或碰撞概率，沿 \cpath{BinaryCollision} 树阅读。两棵树的输入参数、随机变量、守恒检查与适用范围不能互换。

\section{本章结论}

粒子推进器不等于孤立的 Boris 公式。一次 WarpX 粒子更新连接时间层、场插值、动量/位置更新、沉积、边界和多物理创建。面对异常轨道、错误粒子数或不守恒的场，可按三步缩小问题：

\begin{enumerate}
\item \textbf{先固定观察量和时间层。} 位置、机械动量、\cpath{rho}、current、scraped buffer 和 reduced diagnostic 不一定在同一时刻；先写清比较对象。
\item \textbf{再追带电粒子主链。} \cpath{PushParticlesandDeposit()} 经容器 tile loop 依次组织旧 \cpath{rho}、gather、Boris/Vay/Higuera--Cary/RR、\cpath{UpdatePosition()}、current 和新 \cpath{rho}。
\item \textbf{最后拆开独立分支。} particle boundary、PML、AMR 与产生模型需要各自的 observable；强场 QED 沿 \code{chi -> optical depth -> lookup table -> source/product}，virtual-photon/linear 碰撞沿 \cpath{BinaryCollision}，两者不能互换参数或守恒检查。
\end{enumerate}

这也给出了后续章节接口：第 5 章解释守恒沉积，第 6 章解释场推进，第 7 章解释 PML/AMR/边界，第 8 章将上述量写成可判断的 diagnostics。

\section{练习与复现实验}

\begin{enumerate}
\item \textbf{pusher 对照题}：比较 Boris、Vay 与 Higuera-Cary，说明 Boris 的大位移为何需要回到 force-free relativistic pusher 的适用条件，而非简单归为“代码运行失败”。
\item \textbf{源码定位题}：从 \cpath{PhysicalParticleContainer::Evolve()} 追到 gather、momentum push、\cpath{UpdatePosition} 和 current/charge deposition，画出 tile loop 的四个时间层节点。
\item \textbf{最小复现实验}：运行 \cpath{particle\_pusher} 或 \cpath{photon\_pusher} analysis，记录位置、动量和电流沉积；解释带电 pusher 与无质量 photon 为何不能共用位置/动量容差。
\item \textbf{QED 路线题}：选一个 quantum-synchrotron 或 Breit-Wheeler case，列出 source、product、\cpath{opticalDepthQSR/BW}、主观察量和一个不可外推结论，并说明它不能替代 \cpath{linear\_breit\_wheeler} 的碰撞参数。
\end{enumerate}
